\documentclass[11pt]{article}
\usepackage[margin=1in]{geometry}
\usepackage[T1]{fontenc}
\usepackage[utf8]{inputenc}
\usepackage{lmodern}
\usepackage{amsmath,amssymb,amsthm,mathtools}
\usepackage{booktabs,tabularx,longtable,array}
\usepackage{xltabular}% tabularx-style tables that can break across pages
\usepackage{enumitem}
\usepackage[dvipsnames]{xcolor}
\usepackage{cite}% sorts and compresses numeric citations
\usepackage[colorlinks=true,linkcolor=Blue,citecolor=Blue,urlcolor=Blue]{hyperref}

% ---------------------------------------------------------------------
% Shared notation
\newcommand{\R}{\mathbb{R}}
\newcommand{\E}{\mathbb{E}}
\renewcommand{\P}{\mathbb{P}}
\newcommand{\one}{\mathbf{1}}
\newcommand{\D}{\mathcal{D}}
\newcommand{\Lgen}{\mathcal{L}}
\DeclareMathOperator{\Var}{Var}
\DeclareMathOperator{\Cov}{Cov}
\DeclareMathOperator{\Ent}{Ent}
\DeclareMathOperator{\KL}{KL}
\DeclareMathOperator{\Tr}{Tr}
\DeclareMathOperator{\ran}{ran}
\newcommand{\norm}[1]{\lVert #1\rVert}
\newcommand{\ip}[2]{\langle #1,#2\rangle}
% Locators.  \ploc{compiled locator}{TeX file:line} cites the preprint (or a
% companion).  Only the compiled locator (section, equation, result, page of the
% PDF) is printed.  The TeX file:line stays in the source to help checkers find
% the passage and never appears in the PDF.  \sloc{key}{locator} cites a source.
\newcommand{\ploc}[2]{#1}
\newcommand{\sloc}[2]{\cite[#2]{#1}}
% Relationship types and flags, as used in the summary table and the Summary
% paragraphs: \rel{same technique}; \attrgap and \overstated tag a Cited?
% cell (precedence_stats.py checks them against the JSON).
\newcommand{\rel}[1]{\textit{#1}}
\newcommand{\attrgap}{\textbf{attribution gap}}
\newcommand{\overstated}{\textbf{overstated novelty}}
% Uniform list and table styling across all findings.
\setlist{itemsep=2pt,topsep=3pt,parsep=0pt,leftmargin=1.6em}
\setlist[enumerate]{label=\arabic*.}
\setlist[itemize]{label=\textbullet}
\renewcommand{\arraystretch}{1.15}
% Findings are numbered by LaTeX, so nobody types an ID.  In the main file,
% \findinggroup{F} after the Findings \section sets the ID prefix and restarts
% the count; IDs follow the order of the \input lines (by magnitude).  Each
% fragment opens with
%   \finding{<title>}{<slug>}   or   \finding[<short title>]{<title>}{<slug>}
% which prints an unnumbered subsection "F3. <title>", adds "F3. <title>" (or
% the short title) to the table of contents and PDF bookmarks, and does
% \label{f:<slug>}, so \ref{f:<slug>} prints the ID "F3" as a link.  The
% heading carries no LaTeX section number, so a finding never shows two
% numbers (no "3.2 F2").  Refer to another finding as finding~\ref{f:<slug>}.
% Before any \findinggroup the IDs print as ?1, ?2, ... (as in a standalone
% test of one fragment).
\makeatletter
\newcounter{finding}
\newcommand{\finding@prefix}{?}
\renewcommand{\thefinding}{\finding@prefix\arabic{finding}}
\newcommand{\findinggroup}[1]{\gdef\finding@prefix{#1}\setcounter{finding}{0}}
\newcommand{\finding}[3][]{%
  \stepcounter{finding}%
  \protected@edef\finding@id{\thefinding}% fix the ID now: "F3", not a counter reference
  \edef\finding@next{\noexpand\finding@sect{\finding@id}}%
  \finding@next{#1}{#2}{#3}}
% #1 ID (already expanded), #2 short title (may be empty), #3 title, #4 slug.
\newcommand{\finding@sect}[4]{%
  \subsection*{#1. #3}% hyperref anchors \subsection* itself, just above the heading
  \if\relax\detokenize{#2}\relax
    \addcontentsline{toc}{subsection}{#1. #3}%
  \else
    \addcontentsline{toc}{subsection}{#1. #2}%
  \fi
  \def\@currentlabel{#1}% \ref prints the ID; hyperref links it to this heading
  \label{f:#4}}
\makeatother
% ---------------------------------------------------------------------
\space%

\title{Missing citations in\\ ``A spectral gap throughout the high-temperature Sherrington--Kirkpatrick phase''}
\date{October 13, 2026}

\begin{document}
\maketitle

\section*{Summary of findings}

Each row is a citation the preprint is missing: earlier work that already
contains part of the preprint, which the preprint does not cite where that
part appears. ``Cited?'' names the earlier work and says what, if anything, the
preprint cites it for. The exception is a row flagged \emph{overstated
novelty}: there the preprint cites the work, but its framing presents the
content as new. Terms are defined in Section~\ref{sec:about}; centrality is a
best-effort judgment. Locations use the preprint's own numbering and printed
pages.

\begin{center}\footnotesize
\begin{xltabular}{\linewidth}{@{}l >{\raggedright\arraybackslash}p{0.11\linewidth} >{\raggedright\arraybackslash}p{0.10\linewidth} >{\raggedright\arraybackslash}X >{\raggedright\arraybackslash}p{0.17\linewidth} >{\raggedright\arraybackslash}p{0.24\linewidth}@{}}
\toprule
ID & Relationship & Centrality & Technique or result & In the preprint & Cited?\\
\midrule
\endfirsthead
\toprule
ID & Relationship & Centrality & Technique or result & In the preprint & Cited?\\
\midrule
\endhead
\ref{f:planted-law-transfer} & \rel{same technique} & central &
Planted law of disorder and observation path, likelihood ratio $Z/\E Z$, and transfer back to the original disorder &
(5.4) (p.~28); size bias, (5.27) and transfer (pp.~39--40) &
El Alaoui--Montanari--Sellke 2022 \cite{ElAlaouiMontanariSellke2022-SamplingSKAlgorithmicSL}: credited elsewhere (named in \S1.1, p.~4, not in \S5); Aizenman--Lebowitz--Ruelle 1987 \cite{NL-AizenmanLebowitzRuelle1987}: credited elsewhere (p.~3)\\
\addlinespace
\ref{f:stopped-localization} & \rel{adapted} & central &
Stopping the localization before its path estimates can fail, with early stops and terminal failures as one exceptional event &
\S2.2, (2.3)--(2.4) (pp.~7--8); \S2.4 (p.~10); Cor.~5.2 (p.~27) &
Huang--Mohanty--Rajaraman--Wu 2024 \cite{HuangEtAl2024-WeakPoincareSimulatedAnnealing}: not cited; Chen--Eldan 2022 \cite{ChenEldan2022-LocalizationSchemes}: credited in part; \attrgap\\
\addlinespace
\ref{f:uniform-diagonal-equivalents} & \rel{adapted} & central &
Deterministic equivalents for the diagonal of words with an Onsager-shifted inverse, uniform over diagonal parameters chosen after the matrix &
Lemmas~3.1 and 3.3, (3.18) (pp.~12--19) &
Jekel--Sandhu--Shi 2024 \cite{JekelSandhuShi2024-PotentialHessianAscentSK}; Davies--Lee--Sandhu--Shi 2026 (PHA~III) \cite{DaviesEtAl2026-PotentialHessianAscentIIISampling}; Fan--Mei--Montanari 2018 \cite{FanMeiMontanari2018-TAPFreeEnergyVariationalInference}: not cited; Speicher 1998 \cite{Speicher1998-OperatorValuedFreeProbability}: credited in part; \attrgap\\
\addlinespace
\ref{f:planted-tap-kac-rice} & \rel{adapted} & central &
Planted Kac--Rice first moment for TAP critical points: determinant rate, Gaussian density of the TAP map, complexity functional and its sign &
Lemma~3.5 (p.~20); Lemma~4.1 (pp.~22--24); (5.5), Lemma~5.3, (5.9), Lemma~5.4 (pp.~28--31) &
Fan--Mei--Montanari 2018 \cite{FanMeiMontanari2018-TAPFreeEnergyVariationalInference}: not cited; \attrgap\\
\addlinespace
\ref{f:approximate-zero-local-volume} & \rel{adapted} & central &
Control of approximate critical points from an annealed first moment: smoothed integral, local-volume lower bound, determinant cancellation &
(5.5) (p.~28); Lemma~5.6 (pp.~35--37); exclusion (pp.~37--38); Lemma~3.5 (pp.~20--21) &
Huang--Sellke 2023 \cite{HuangSellke2023-StrongTopologicalTrivialization}, Huang--Montanari--Pham 2024 \cite{HuangMontanariPham2024-SphericalSpinGlassSamplingSL}: not cited; \attrgap\\
\addlinespace
\ref{f:tap-stability-route} & \rel{adapted} & central &
Well-posedness of the TAP equation along the localization path: planted law, first-moment exclusion with a conditional Hessian, approximate critical points, transfer &
Prop.~5.1 (p.~27); \S5.4 (pp.~37--40); Lemma~5.7 (pp.~40--41) &
Huang--Montanari--Pham 2024 \cite{HuangMontanariPham2024-SphericalSpinGlassSamplingSL}: not cited; \attrgap\\
\addlinespace
\ref{f:conditional-hessian-first-moment} & \rel{adapted} & central &
TAP Hessian conditioned on a vanishing gradient, positive near the Bayes law, inside a Kac--Rice first moment &
(5.8) (pp.~28--29); Lemma~5.5 (pp.~31--35); (5.25) (p.~37) &
Celentano--Fan--Mei 2021 \cite{CelentanoFanMei2021-LocalConvexityTAPZ2Sync}: not cited; \attrgap\\
\addlinespace
\ref{f:amp-memory-onsager-rule} & \rel{adapted} & central &
A memory correction for every input of a multi-input recursion, and its cancellation of one-step reversals at every depth &
Def.~6.2, (6.7) (p.~44); (6.14) (p.~46); proof of Lemma~6.5 (p.~54); proofs of Lemmas~6.1, 6.4 and 3.3 (pp.~43, 51--52, 16) &
Bayati--Lelarge--Montanari 2012 \cite{BayatiLelargeMontanari2012-UniversalityMessagePassing}; Montanari 2018 \cite{Montanari2018-OptimizationSKHamiltonian}: not cited (Thouless--Anderson--Palmer cited for the Onsager correction, Bolthausen for (6.1); AMP named only in \S1.1, p.~4, citing others for other content); \attrgap\\
\addlinespace
\ref{f:terminal-gap-sparse-submatrices} & \rel{adapted} & central &
Gap at a large observation time from few weak-field coordinates and uniformly small principal submatrices &
Lemma~8.2, Prop.~8.3 (pp.~61--64) &
Davies--Lee--Sandhu--Shi 2026 (PHA~III) \cite{DaviesEtAl2026-PotentialHessianAscentIIISampling}, Davies--Lee--Sandhu--Shi 2026 (weak Poincar\'e) \cite{DaviesEtAl2026-WeakPoincareApproxSL}: not cited; Bandeira--El Alaoui--R\"odder 2026 \cite{BandeiraElAlaouiRodder2026-GlauberHighOverlapGibbs}: credited in part; \attrgap\\
\addlinespace
\ref{f:tap-hessian-covariance-route} & \rel{same route} & central &
Poincar\'e inequality for all $\beta<1$ from operator-norm covariance bounds at the localization fields, via a uniformly positive TAP Hessian &
\S1.4 (p.~6); Props.~5.1, 7.2; \S2 &
El Alaoui--Gaitonde 2022 \cite{ElAlaouiGaitonde2022-CovarianceMatrixSK}: cited for other content (Thm~1.1); Brennecke--Xu--Yau 2023 \cite{BrenneckeXuYau2023-OperatorNormCorrelationSK}: not cited\\
\addlinespace
\ref{f:weak-poincare-beta-half} & \rel{weaker result} & central &
Weak Poincar\'e inequality for zero-field SK heat-bath dynamics at $\beta<1/2$ &
Abstract; Thm~1.1 (p.~3) &
Davies--Lee--Sandhu--Shi 2026 (weak Poincar\'e) \cite{DaviesEtAl2026-WeakPoincareApproxSL}: not cited\\
\addlinespace
\ref{f:sudakov-fernique-tap-margin} & \rel{same technique} & supporting &
Margin for the Onsager-shifted TAP Hessian by comparing $v^{\mathsf T}Wv$ with a norm-weighted linear process under a common shift &
Lemma~3.2, (3.8)--(3.9) (pp.~12--13) &
Celentano--Fan--Mei 2021 \cite{CelentanoFanMei2021-LocalConvexityTAPZ2Sync}: not cited; Celentano 2022 \cite{Celentano2022-SudakovFerniquePostAMPTAPConvexity}: cited for other content; \attrgap\\
\addlinespace
\ref{f:nishimori-consistent-gaussian} & \rel{same technique} & supporting &
Consistent Gaussian law: mean equals variance on the Nishimori line, then a one-variable fixed point with uniqueness, size and Lipschitz bounds &
Lemma~4.2, (4.7)--(4.11) (pp.~25--26) &
El Alaoui--Montanari--Sellke 2022 \cite{ElAlaouiMontanariSellke2022-SamplingSKAlgorithmicSL}: cited for other content; Davies--Lee--Sandhu--Shi 2026 (PHA~III) \cite{DaviesEtAl2026-PotentialHessianAscentIIISampling}: not cited; \attrgap\\
\addlinespace
\ref{f:tilted-path-entropy} & \rel{fragment} & supporting &
Path-law entropy of the reweighted observations as the entropy lost along the localization &
\S2.3, (2.5) and (2.7) (pp.~8--9) &
Chen--Eldan 2022 \cite{ChenEldan2022-LocalizationSchemes}: cited for other content (Lehec \cite{Lehec2010-RepresentationFormulaEntropy} is cited for the identity)\\
\addlinespace
\ref{f:mixture-gaussian-stein-kernel} & \rel{fragment} & supporting &
Stein operator of the cosh-tilted Gaussian mixture and its bounded solution for $\tanh$ &
(6.8)--(6.9) (p.~45), used in Lemma~6.5 (pp.~53--54); (6.12)--(6.13) (pp.~45--46); (7.3) (p.~57) &
Chatterjee 2007 \cite{Chatterjee2007-SpinGlassesSteinsMethod}: not cited\\
\addlinespace
\ref{f:posterior-mean-tap-root} & \rel{fragment} & supporting &
Comparing a posterior mean with the TAP root through a late Onsager-corrected iterate &
(6.10)--(6.11) (p.~45); Prop.~7.3 (p.~58) &
El Alaoui--Montanari--Sellke 2022 \cite{ElAlaouiMontanariSellke2022-SamplingSKAlgorithmicSL}: credited elsewhere (\S1.1, p.~4); Celentano--Fan--Mei 2021 \cite{CelentanoFanMei2021-LocalConvexityTAPZ2Sync}: not cited\\
\bottomrule
\end{xltabular}
\end{center}

\clearpage
\space%

\tableofcontents

\section{About this document}\label{sec:about}

\paragraph{The preprint.}
The preprint is \emph{A spectral gap throughout the high-temperature
Sherrington--Kirkpatrick phase}, dated September 24, 2026, with author line
``OpenAI'' \cite{OAI-Preprint}. It is one of the manuscripts of the openai/math
release. It proves that for every fixed inverse temperature $0<\beta<1$, with
probability tending to one over the disorder, the zero-field Gaussian
Sherrington--Kirkpatrick Gibbs measure satisfies a Poincar\'e inequality for
the unscaled single-site heat-bath Dirichlet form, with a constant depending
only on $\beta$; equivalently, the heat-bath chain has spectral gap of order
$1/n$. Section~\ref{sec:overview} summarizes the paper and fixes the notation
used throughout this document. The preprint cites no other preprint from the
same release.

\paragraph{What this document reports.}
Each finding is a citation the preprint is missing: part of the preprint's
mathematics (a technique, construction, estimate, proof route or result) that
an identifiable earlier paper already contains, where the preprint does not
cite that paper for it where it is used. ``Earlier'' means publicly available
before the preprint's date. The one exception, overstated novelty, is
described below. Not reported: textbook material in its textbook form; weaker versions of
intermediate statements; and, for content the preprint uses, cases where it
cites a later source that covers that content rather than the original. An
earlier result on the same problem is reported even when later work that the
preprint cites implies it. Findings are ordered by centrality, then by relationship in
the order listed below, then by position in the preprint.

\paragraph{Relationships.}
A finding's relationship is the one between the preprint and the first
earlier work listed for it (the principal), the work with the highest relationship in the
order below; when several hold that relationship, however closely each matches,
the one the others credit or build on, otherwise the earliest. Any other works listed relate to the same content more weakly or
in part; the finding says how, including which work originated the idea when
that is not the first.
\begin{description}[leftmargin=1.5em,font=\normalfont\itshape]
\item[same result.] An earlier paper, or several earlier papers together, prove essentially the same statement.
\item[same technique.] The earlier technique, construction or theory is
reused essentially intact, possibly in a new setting or with new ingredients
around it.
\item[adapted.] The earlier technique's structure is followed, with a change
at its core: a different object, or a key step replaced.
\item[proof shortcut.] A published result would replace or substantially
shorten a large part of the argument.
\item[same route.] An earlier paper proposed the proof strategy, even
heuristically, and the preprint carries it out.
\item[fragment.] One identifiable piece of the earlier technique is reused.
\item[weaker result.] An earlier paper proves a weaker, special-case or
partial-range version of one of the preprint's main results, or of a
statement it presents as a significant contribution.
\end{description}

\paragraph{Centrality.}
Each finding also gets a best-effort judgment of how much of the preprint's
main argument depends on it: \emph{main result} (reserved for earlier work that
covers the main result itself, or most of it), \emph{central},
\emph{supporting} or \emph{peripheral}. An earlier weaker result is never
labelled main result. Each finding states the chain of results it feeds, so
the basis for the judgment is visible.

\paragraph{Citations.}
The ``Cited?'' column of the summary names the earlier work and says what the
preprint cites for the content in question: \emph{not cited} (the work is not
in the bibliography, and nothing the preprint cites credits it for this
content), \emph{cited for other content}, \emph{credited elsewhere} (credited
for this content, directly or through a cited location that credits it, but
not where it is used, in a way that an expert reader who doesn't know the earlier work might not connect with it), or \emph{credited in part}. A citation counts only when it is \emph{for that specific content};
citing the same paper nearby for a different result does not count. Content
the preprint credits where it is used, directly or through a cited location
that itself credits the earlier work, is not reported, except when its novelty is overstated (below). For
overstated novelty, the column says where the preprint does cite the work.

\paragraph{Flags.}
\begin{description}[leftmargin=1.5em,font=\normalfont\bfseries]
\item[Attribution gap.] Marks the most serious missing citations: the
relationship is same technique, adapted or same result; the centrality is main
result, central or supporting; and the preprint cites the principal earlier
work for this content nowhere.
\item[Overstated novelty.] The preprint does cite the earlier work for this
content, but its abstract, a statement of its significant contributions, or the framing of
a main result presents that content as new, explicitly or in the way a reader
would reasonably take it. This is the one case in which content the preprint
credits is reported.
\end{description}

\paragraph{How to check a claim.}
Locations in the preprint use its section,
equation and theorem numbers and the page numbers printed in its PDF.
Locations in other sources give the theorem, equation and printed page of the
version listed in the bibliography. Each finding ends with a short checklist
of the passages to put side by side.

\paragraph{Provenance.}
This document was generated by LLM agents (Claude Opus 5.5) running the
precedence-review skill
(\href{https://github.com/jnshi/math-lineage/tree/e025345ea/.agents/skills/precedence-review}{jnshi/math-lineage},
commit \texttt{e025345ea}), makes no claim of completeness, and has not been
vetted by a human. The run was used to develop the skill, and its later stages
ran revisions of the skill made during the run, the last of them commit
\href{https://github.com/jnshi/math-lineage/tree/e1e8ffbf7/.agents/skills/precedence-review}{\texttt{e1e8ffbf7}}.
It was run by Jonathan Shi. The leads for this run came from the findings of
an earlier run with version 1 of the skill, in which Jonathan Shi supplied the
lead that became finding~\ref{f:uniform-diagonal-equivalents}. Jonathan Shi
makes no claim of impartiality.
\space%
\section{The preprint in brief, and notation}\label{sec:overview}

\paragraph{Model and theorem.}
Fix $0<\beta<1$ and put $j=\beta^2$. The disorder $J$ is a real symmetric
$n\times n$ matrix with zero diagonal and independent entries
$J_{ij}\sim N(0,j/n)$, $i<j$. For a field $h\in\R^n$,
\[
 \mu_h(x)=\frac{1}{Z_h}\exp\Bigl\{\tfrac12x^{\mathsf T}Jx+h^{\mathsf T}x\Bigr\},
 \qquad x\in\{-1,1\}^n,
\]
and $\mu_0$ is the zero-field Gibbs measure. With $P_i^h$ the conditional
expectation under $\mu_h$ given all spins but the $i$th, the unscaled
heat-bath Dirichlet form is
$\D_h(f)=\sum_{i}\E_{\mu_h}(f-P_i^hf)^2$ \ploc{(1.1), p.~3}{introduction.tex:22--24}.
The main result, \ploc{Theorem~1.1 (p.~3)}{introduction.tex:33--42}, states
that there is $C_\beta<\infty$ with
\[
 \P_J\bigl\{\Var_{\mu_0}(f)\le C_\beta\,\D_0(f)\text{ for every }f:\{-1,1\}^n\to\R\bigr\}\longrightarrow1,
\]
so the discrete-time heat-bath chain that updates a uniform site has
spectral gap at least $1/(C_\beta n)$ with probability tending to one.

\paragraph{The proof, in the preprint's own terms.}
The preprint describes two parts: ``prove useful inequalities for posterior
Gibbs laws, then transfer them to the unobserved law''
\ploc{(\S1.3, p.~5)}{introduction.tex:173--174}. It observes a spin sample in
Gaussian noise,
\[
 X\sim\mu_0,\qquad Y(t)=tX+B(t),\quad 0\le t\le T,
\]
with $B$ a standard Brownian motion, as in \ploc{(1.4), p.~6}{introduction.tex:199--201},
so that the posterior spin law at time $t$ is $\mu_{Y(t)}$. The transfer is
proved in the preprint's \S2 from three inputs
\ploc{(\S1.4, p.~6)}{introduction.tex:197--237}:
\begin{enumerate}
\item stability of every approximate solution of the field equation along
ordinary observation paths, except on an event of probability
$e^{-an}$ \ploc{(Proposition~5.1, p.~27)}{stability.tex:35--60}, proved in
\S5 with random-matrix estimates from \S3 and a scalar entropy inequality
from \S4;
\item two estimates at such a good field $h$, proved in \S7 from the residual
identities of \S6: the directional covariance bound
\[
 \norm{\Cov_{\mu_h}(G,x)}^2\le C\bigl(\Var_{\mu_h}(G)+\D_h(G)\bigr)
\]
of \ploc{(1.5), p.~6}{introduction.tex:214--217} and
\ploc{Proposition~7.2, p.~56}{consequences.tex:89--98}, and the comparison of
the mean of $\nu=f^2\mu_h/\E_{\mu_h}f^2$ with the root in
\ploc{(1.6), p.~6}{introduction.tex:218--224} and
\ploc{Proposition~7.3, p.~58}{consequences.tex:202--213};
\item a terminal gap $\Var_{\mu_{Y(T)}}f\le2\D_{Y(T)}(f)$ at a large fixed
time $T$ \ploc{(Proposition~8.3, p.~63)}{terminal.tex:334--433}, proved in \S8.
\end{enumerate}

\paragraph{Field equation and stability.}
For $y\in\R^n$ put $m(y)=\tanh y$, $q(y)=n^{-1}\sum_i m_i(y)^2$,
$b(y)=1-q(y)$ and $V_y=\operatorname{diag}(1-m(y)^2)$. The field map of
\ploc{(5.1), p.~26}{stability.tex:9--20} is
\[
 F_{J,h_0}(y)=y-h_0+\bigl(jb(y)I-J\bigr)m(y),
\]
and for a nonnegative diagonal matrix $A$ the stability matrix of
\ploc{(5.2), p.~27}{stability.tex:22--28} is
\[
 \mathsf H_J(y,A)=I+A^{1/2}\Bigl(jb(y)I-J-\frac{2j}{n}m(y)m(y)^{\mathsf T}\Bigr)A^{1/2},
\]
the symmetric version of the derivative $F'_{J,h_0}(y)$ when $A=V_y$.
At the fields under consideration, $r(h)$ denotes the unique zero of
$F_{J,h}$ and $t_*(h)=\tanh r(h)$
\ploc{(\S1.4, p.~6)}{introduction.tex:225--231}.

\paragraph{Matrix notation of \S3.}
The preprint couples $J$ to a centered GOE matrix $W$ with
$\E W_{ab}^2=j/n$ ($a<b$) and $\E W_{aa}^2=2j/n$ through
$J=W-\operatorname{diag}(W)$. For a nonnegative diagonal $A$ with entries at
most $1+\eta$, where $j(1+\eta)^2<1$, it sets $\chi_A=n^{-1}\Tr A$ and, in
\ploc{(3.2), p.~10}{matrices.tex:26--37},
\[
 K_A(z;M)=(I-zAM+z^2j\chi_AA)^{-1},\qquad
 S_A(z;M)=I+z^2j\chi_AA-zA^{1/2}MA^{1/2}.
\]

\paragraph{The recursion of \S6.}
Started from a Gibbs sample $x$ at a field $h_0$, the primary recursion of
\ploc{(6.1), p.~42}{residuals.tex:36--43} is
\[
 m^0=x,\quad b_0=0,\quad m^l=\tanh h^l,\quad
 b_l=\frac1n\sum_i\bigl(1-(m^l_i)^2\bigr),\quad
 h^{l+1}=h_0+Jm^l-jb_lm^{l-1},
\]
with $h^1=h_0+Jx$ and residuals $R_l=m^{l-1}-m^l$.

\space%

% IDs are assigned by LaTeX in the order of the \input lines (F1, F2, ...).
% Order by magnitude: centrality (main result, central, supporting,
% peripheral), then the relationship (same result, same technique, adapted,
% proof shortcut, same route, fragment, weaker result), then reading order.
% precedence_stats.py checks this order.
% Withdrawing or reordering a finding is just deleting or moving its \input line.
\section{Findings}\label{sec:findings}
\findinggroup{F}
\finding[Planted law and partition-function transfer]{A planted law for the disorder and the observation path, transferred back through the normalized partition function}{planted-law-transfer}

\paragraph{Summary.}
El Alaoui, Montanari and Sellke analyse their Sherrington--Kirkpatrick
sampler under a planted law of disorder and observation path, prove that its
likelihood ratio against the original law is the normalized partition
function, and transfer back by contiguity
\sloc{ElAlaouiMontanariSellke2022-SamplingSKAlgorithmicSL}{\S4.1, Prop.~4.2, pp.~10--12}.
The preprint names that paper only in its introduction
(\ploc{\S1.1, p.~4}{introduction.tex:118--120}), for combining ``Gaussian
observations, planted laws, and approximate message passing''; in its \S5,
where it introduces its planted law \ploc{(5.4), p.~28}{stability.tex:107--118}
and derives the size bias and the transfer back
(\ploc{pp.~39--40}{stability.tex:1022--1083}), it cites nothing.
The shared approach: the computations along the path need an explicit law,
and size-biasing the law of disorder, spin sample and path by the normalized
partition function gives one (uniform spins, a GOE matrix plus a
rank-one spike, path $tx+B$); as that ratio is not exponentially small for
$\beta<1$, failure probabilities that vanish under the planted law
(exponentially fast, in the preprint's form) also vanish under the original
one, an insight that, in this setting, is El Alaoui, Montanari and Sellke's.
\rel{same technique}: (5.4) is their planted law conjugated by
$D_{x_0}=\operatorname{diag}(x_0)$, the size bias on p.~39 is their
Proposition~4.2 with the GOE diagonal split off, and the Markov-inequality
transfer is a quantitative form of their contiguity step.
The lower tail \ploc{(5.27)}{stability.tex:1040--1065}, proved anew, follows
from the Aizenman--Lebowitz--Ruelle limit \cite{NL-AizenmanLebowitzRuelle1987}
quoted on \ploc{p.~3}{introduction.tex:53--56}, credited elsewhere.

\paragraph{Correspondence.}
\begin{center}\small
\begin{xltabular}{\linewidth}{@{}>{\raggedright\arraybackslash}p{0.09\linewidth}>{\raggedright\arraybackslash}X>{\raggedright\arraybackslash}X@{}}
\toprule
 & Preprint & El Alaoui--Montanari--Sellke\\
\midrule
\endfirsthead
\toprule
 & Preprint & El Alaoui--Montanari--Sellke\\
\midrule
\endhead
Law & $(J,X)$, $X\sim\mu_0$, biased by $L(J)=Z(J)/\E Z(J)$, $Z(J)=Z_0$: $X$ uniform, $J_{ik}\mid X$ independent, mean $jX_iX_k/n$ (\ploc{p.~39}{stability.tex:1022--1033}) & $\mu_{\mathrm{pl}}$: $x$ uniform, $A=(\beta/n)xx^{\mathsf T}+W$, $W\sim\mathrm{GOE}(n)$; $A$-marginal $\mu_{\mathrm{GOE}}(dA)\,Z_{\mathsf{SK}}(A)$ ((4.1)--(4.5), p.~10)\\
\addlinespace
Path & $Y(t)=tX+B(t)$ \ploc{(1.4)}{introduction.tex:199--201}; its law given $J$ is unchanged by the size bias (\ploc{p.~39}{stability.tex:1037--1038}) & $\mathbb Q$: $A\sim\mathrm{GOE}(n)$, $y\overset{d}{=}tx+\bar B$, $x\sim\mu_A$ (4.12); $\mathbb P$: $y=tx_0+B$ (4.9); $d\mathbb P/d\mathbb Q=Z_{\mathsf{SK}}(A)$ (Prop.~4.2, p.~11)\\
\addlinespace
Planted model & $\tilde J=W+(j/n)\one\one^{\mathsf T}$, $h_0=t\one+\sqrt t\,g$ \ploc{(5.4), p.~28}{stability.tex:107--118} & $\mathbb P$: $\beta A=(j/n)x_0x_0^{\mathsf T}+\beta W$, $y(t)=tx_0+B(t)$; their proof works under $\mathbb P$ (p.~12)\\
\addlinespace
Transfer & $\P\{L(J)<e^{-\delta n}\}\to0$ \ploc{(5.27)}{stability.tex:1040--1065}; with Markov, the planted bound $Ce^{-a_4n}$ gives path failure $\le e^{-an}$ given typical $J$ (\ploc{pp.~39--40}{stability.tex:1067--1083}) & $Z_{\mathsf{SK}}(A)\to e^{W}$ in law, $W\sim N(-\sigma^2,2\sigma^2)$ (Thm~4.1, p.~11); by Le Cam's first lemma, results under $\mathbb P$ transfer to $\mathbb Q$ (pp.~11--12)\\
\addlinespace
\bottomrule
\end{xltabular}
\end{center}

\paragraph{How to check.}
\begin{enumerate}
\item The size bias on \ploc{p.~39}{stability.tex:1022--1038} against
\sloc{ElAlaouiMontanariSellke2022-SamplingSKAlgorithmicSL}{(4.1)--(4.6), p.~10; Prop.~4.2, p.~11}.
For $J=\beta(A-\operatorname{diag}A)$, $\mu_A=\mu_0$ and
$Z_{\mathsf{SK}}(A)=2^{-n}\sum_xe^{\beta\ip{x}{Ax}/2-\beta^2n/4}=L(J)\,e^{\beta\Tr A/2-j/4}$,
as $\E Z(J)=2^ne^{j(n-1)/4}$; $\beta\Tr A/2\sim N(0,j/2)$ is independent of
$J$, and tilting by $e^{\beta\Tr A/2-j/4}$ gives each $\beta A_{ii}$ the law
$N(j/n,2j/n)$ that the preprint adds by hand.
\item \ploc{(5.4), p.~28}{stability.tex:107--118} against
\sloc{ElAlaouiMontanariSellke2022-SamplingSKAlgorithmicSL}{(4.3), p.~10; (4.9), p.~11}:
$\tilde J=D_{x_0}(\beta A)D_{x_0}$ and $h_0=D_{x_0}y(t)$.
\item The transfer on \ploc{pp.~39--40}{stability.tex:1040--1083} against
\sloc{ElAlaouiMontanariSellke2022-SamplingSKAlgorithmicSL}{Thm~4.1, p.~11; p.~12};
(5.27) against \ploc{p.~3}{introduction.tex:53--56}: with $\E Z(J)$ as in
item~1, the limit quoted there gives $n^{-1}\log L(J)\to0$ in probability,
hence (5.27).
\end{enumerate}

\space%
\finding[Stopped localization with one exceptional event]{Stopping the localization before its path estimates can fail, with early stops and terminal failures as one exceptional event}{stopped-localization}

\paragraph{Summary.}
Huang, Mohanty, Rajaraman and Wu introduce ``stopped localization schemes''
as a generic tool \sloc{HuangEtAl2024-WeakPoincareSimulatedAnnealing}{\S6, pp.~25--27},
and the preprint does not cite their paper; it credits Chen and Eldan for the
variance-retention and Dirichlet-form framework
(\ploc{\S1.1, p.~4}{introduction.tex:105--110}; \ploc{\S2, p.~7}{observation.tex:7--9}),
whose annealing theorem needs the terminal gap almost surely
\sloc{ChenEldan2022-LocalizationSchemes}{Thm~46, p.~32}.
The approach: when the estimate driving variance retention and the terminal
inequality hold only on typical paths, stop where the estimate may
fail, so that retention holds by construction (their Lemma~6.6, p.~26), and
treat early stops and terminal failures as one exceptional event
(Lemma~6.8, p.~27); the combination is their insight.
The preprint's \ploc{\S2.2 (pp.~7--8)}{observation.tex:45--88} and
\ploc{\S2.4 (p.~10)}{observation.tex:188--205} follow it, stopping where the
stability of \ploc{Corollary~5.2 (p.~27)}{stability.tex:62--97} may fail:
\rel{adapted}, since the exceptional event is charged by its probability under
the path law started from $f^2\mu_0$, not by $\mathrm{osc}(f)^2$ times its
probability, which gives a Poincar\'e inequality instead of a weak one.
The moves are older (fragments, not cited for this): Eldan stopped stochastic
localization when good events failed
\sloc{Eldan2012-ThinShellSpectralGapStochLoc}{\S5, p.~26}, and El Alaoui,
Eldan, Gheissari and Piana split the terminal variance into good and bad paths
\sloc{ElAlaouiEtAl2023-FastRelaxationRFIM}{p.~15}; Davies, Lee, Sandhu and
Shi's weak Poincar\'e paper restates Lemma~6.8 on the hypercube
\sloc{DaviesEtAl2026-WeakPoincareApproxSL}{Thm~5.5, p.~22}.

\paragraph{Correspondence.}
Preprint: $\E_{\mu_0}f=0$, $\E_{\mu_0}f^2=1$, $V_t=\Var_{\mu_{Y(t)}}f$,
$N_t=\E_{\mu_{Y(t)}}f^2$, $d_0=\D_0(f)$; $\P^f$, the path law from
$X\sim f^2\mu_0$, has density $N_\tau$ at $\tau$. Huang, Mohanty, Rajaraman
and Wu: $\mu_t$
\sloc{HuangEtAl2024-WeakPoincareSimulatedAnnealing}{Def.~6.2, p.~25} is
$\mu_{Y(t)}$ if $C_t=I$; $(\rho,\delta)$-weak Poincar\'e:
$\Var\varphi\le\rho^{-1}\mathcal E(\varphi,\varphi)+\delta\,\mathrm{osc}(\varphi)^2$.
\begin{center}\small
\begin{xltabular}{\linewidth}{@{}>{\raggedright\arraybackslash}p{0.1\linewidth}>{\raggedright\arraybackslash}X>{\raggedright\arraybackslash}X@{}}
\toprule
 & Preprint & Huang--Mohanty--Rajaraman--Wu\\
\midrule
\endfirsthead
\toprule
 & Preprint & Huang--Mohanty--Rajaraman--Wu\\
\midrule
\endhead
Stopping & $\tau=T\wedge\tau_0$, $\tau_0$ the first time a fixed witness $(y,A)$, $\mathsf H_J(y,A)\not\succ cI$, has $\norm{F_{J,Y(t)}(y)}\le\rho_0\sqrt n$ (\ploc{\S2.2, p.~7}{observation.tex:46--51}; proof of \ploc{Cor.~5.2, p.~27}{stability.tex:74--97}) & $\tau=T\wedge\inf\{t:\norm{C_t^{1/2}\Cov(\mu_t)C_t^{1/2}}_{\mathrm{op}}\ge K\}$ (Lemma~6.6, p.~26)\\
\addlinespace
Retention & $\E V_\tau\ge v-(1-v)d_0$, $v=e^{-C_{\mathrm d}T}$ \ploc{(2.4), p.~8}{observation.tex:70--86} & $\E\Var_{\mu_\tau}[\varphi]\ge e^{-TK}\Var_\mu[\varphi]$ (Lemma~6.6)\\
\addlinespace
Bad event & $\{\tau_0\le T\}\cup\{\tau_0>T,\ Y(T)\notin\mathcal T_n(J)\}$, $\mathcal T_n(J)=\{h:\Var_{\mu_h}g\le2\D_h(g)\ \forall g\}$; $\P(\mathrm{bad})\le2e^{-an}$ \ploc{(2.3), p.~7}{observation.tex:53--67} & covariance bound fails before $T$ (probability ${}\le\eta_1$) or $\mu_T$ has no $(\rho,\delta)$-weak Poincar\'e inequality (${}\le\eta_2$) (Lemma~6.8, p.~27)\\
\addlinespace
Split & good: $\E[\one_{\mathrm{good}}V_\tau]\le2d_0$ by the terminal gap; bad: $\E[\one_{\mathrm{bad}}V_\tau]\le\E[\one_{\mathrm{bad}}N_\tau]=\P^f(\mathrm{bad})\le v/4$ \ploc{(2.9)}{observation.tex:158--160}; so $\Var_{\mu_0}f\le4e^{C_{\mathrm d}T}\D_0(f)$ \ploc{(\S2.4, p.~10)}{observation.tex:188--209} & good: terminal inequality, $\mathcal E_\mu\ge\E\,\mathcal E_{\mu_\tau}$; bad: $\Var_{\mu_\tau}\varphi\le\mathrm{osc}(\varphi)^2$, adding $(\eta_1+\eta_2)\,\mathrm{osc}(\varphi)^2$ (Lemma~4.11, p.~20); so a $(\rho e^{-TK},e^{TK}(\delta+\eta_1+\eta_2))$-weak Poincar\'e inequality\\
\addlinespace
\bottomrule
\end{xltabular}
\end{center}

\paragraph{How to check.}
\begin{enumerate}
\item $\tau$ and (2.4) in \ploc{\S2.2, pp.~7--8}{observation.tex:46--86}
against \sloc{HuangEtAl2024-WeakPoincareSimulatedAnnealing}{Lemma~6.6, p.~26}
with $C_t=I$.
\item \ploc{(2.3), p.~7}{observation.tex:53--67} and the split in
\ploc{\S2.4, p.~10}{observation.tex:188--205} against the proofs of
\sloc{HuangEtAl2024-WeakPoincareSimulatedAnnealing}{Lemma~6.8, p.~27} and
Lemma~4.11 (p.~20): only the bad-path term differs.
\item \ploc{\S1.1, p.~4}{introduction.tex:105--113} against assumption~2 of
\sloc{ChenEldan2022-LocalizationSchemes}{Thm~46, p.~32}: a gap almost surely
at the stopping time.
\end{enumerate}

\space%
\finding{Uniform diagonal equivalents for Onsager-shifted GOE inverses}{uniform-diagonal-equivalents}

\paragraph{Summary.}
Jekel, Sandhu and Shi
\sloc{JekelSandhuShi2024-PotentialHessianAscentSK}{Thm~4.1(3), pp.~30--31}
estimate the diagonal of the Onsager-shifted resolvent of a GOE matrix plus a
diagonal for all diagonals at once, as their algorithm computes the diagonals
from the same matrix; Davies, Lee, Sandhu and Shi
\sloc{DaviesEtAl2026-PotentialHessianAscentIIISampling}{Cor.~B.35, p.~179}
build on this for the squared resolvent. The preprint cites neither, nor Fan,
Mei and Montanari's trace form
\sloc{FanMeiMontanari2018-TAPFreeEnergyVariationalInference}{Lemma~4.9, pp.~19--20}
of the Onsager-point identity (a fragment).
Jekel, Sandhu and Shi's approach, ``uniform asymptotic freeness from the
diagonal'' \sloc{JekelSandhuShi2024-PotentialHessianAscentSK}{p.~30}, predicts
the diagonal by freeness, proves it in expectation by integration by parts and
a Poincar\'e bound, and extends it to all diagonals by concentration over a
net, a step they call standard.
The preprint's \ploc{Lemmas~3.1 and~3.3 (pp.~12--19)}{matrices.tex:114--648}
have the same three-part structure; the relationship is \rel{adapted}, because
they predict by Speicher's operator-valued moment rule, which the preprint
cites \ploc{(p.~11)}{matrices.tex:77--79}, and compute expectations from
finite-$n$ loop equations, not by free-model interpolation.
Davies et~al.\ (also \rel{adapted}) truncate and chain as the preprint does,
and their Lemma~B.18(b) is the free-model form of the identity
$\mathsf P(K_A)=I$ that the preprint proves without citation
\ploc{(p.~16)}{matrices.tex:438--459}: Speicher's functional equation
\sloc{Speicher1998-OperatorValuedFreeProbability}{Thm~4.1.12, pp.~52--53} at
an explicit solution.

\paragraph{Correspondence.}
Here $D=A^{-1}$, so $K_A(J)A=(D+j\chi_AI-J)^{-1}$; the sources'
$2\beta A_{\mathrm{sym}}$ and $\beta A$ have the law of $W$ for
$j=2\beta^2,\beta^2$; $E_{\mathcal D_n}$ is the diagonal part, $S$ is free
semicircular, $b=\beta n^{-\delta}$.
\begin{center}\small
\begin{xltabular}{\linewidth}{@{}>{\raggedright\arraybackslash}p{0.1\linewidth}>{\raggedright\arraybackslash}X>{\raggedright\arraybackslash}X>{\raggedright\arraybackslash}X@{}}
\toprule
 & Preprint & Jekel--Sandhu--Shi \cite{JekelSandhuShi2024-PotentialHessianAscentSK} & Davies--Lee--Sandhu--Shi \cite{DaviesEtAl2026-PotentialHessianAscentIIISampling}\\
\midrule
\endfirsthead
\toprule
 & Preprint & Jekel--Sandhu--Shi \cite{JekelSandhuShi2024-PotentialHessianAscentSK} & Davies--Lee--Sandhu--Shi \cite{DaviesEtAl2026-PotentialHessianAscentIIISampling}\\
\midrule
\endhead
On one event & \ploc{Lemma~3.1 (p.~12)}{matrices.tex:114--134}: for all
$0\le A\le a_+I$ with $S_A(1;J)\succeq cI$ and words $Q$, which may contain
$K_A(J)A$, $\norm{\operatorname{diag}(Q-\mathsf P(Q))}_2\le C$ (Euclidean
norm) &
Thm~4.1(3) (pp.~30--31): for all $D\ge0$ with
$2\beta^2\operatorname{tr}_n(D^{-2})=1$,
$\norm{E_{\mathcal D_n}[P^2]-b(b^2+D^2)^{-1}}_2\le C_3\beta n^{-2\delta}$,
$P^2=-\operatorname{Im}(\tilde z-2\beta A_{\mathrm{sym}}+D)^{-1}$ &
Cor.~B.35 (p.~179): for all $D\succeq I$, $a\in[0,1]$,
$\norm{Y-F'(a)D(F(a)+D)^{-2}D}_{L^2}\le C\beta\gamma^{-4}n^{-1/2}$,
$Y=E_{\mathcal D_n}[D(a-\beta A+D)^{-2}D]$\\
\addlinespace
Onsager point & $\mathsf P(K_A)=I$ \ploc{(p.~16)}{matrices.tex:438--459};
\ploc{(3.18)}{matrices.tex:433--436} at $z=1$: $\E K_{ii}=1+O(n^{-1})$ &
$E_{\mathcal D_n}[(\tilde z-\sqrt2\beta S+D)^{-1}]=(ib+D)^{-1}$,
$\tilde z=ib+2\beta^2\operatorname{tr}_n((ib+D)^{-1})$ (p.~38; Thm~4.7,
Prop.~4.11) &
$E_{\mathcal D_n}[(a-\beta S+D)^{-1}]=D^{-1}$,
$a=\beta^2\operatorname{tr}_n(D^{-1})$ (p.~17 display with Lemma~B.18(b),
p.~164)\\
\addlinespace
Mean & Integration by parts \ploc{(3.15)}{matrices.tex:356--359}, covariance
$\le Cn^{-3/2}$ by \ploc{(3.13)}{matrices.tex:316--319}, loop equations
\ploc{(3.17)}{matrices.tex:407--410} &
Free-model interpolation, integration by parts, Poincar\'e (Prop.~4.20,
Lemma~4.24, pp.~43--47) &
The same, with a cutoff (Prop.~B.23, pp.~169--170)\\
\addlinespace
All diagonals & Projection $\Pi_R$ onto $\{\norm M\le R\}$
\ploc{(p.~13)}{matrices.tex:247--249}; dyadic meshes
\ploc{(3.21)--(3.22), p.~18}{matrices.tex:555--610} &
Herbst concentration, $\varepsilon$-net in $D^{-1}$ (Lemmas~4.18--4.19,
pp.~40--43) &
Projection $\Pi_{C_{\mathrm{op}}}$ onto $\{\norm M_{\mathrm{op}}\le
C_{\mathrm{op}}\}$, chaining (Thm~B.9, pp.~160--162)\\
\addlinespace
\bottomrule
\end{xltabular}
\end{center}

\paragraph{How to check.}
\begin{enumerate}
\item \sloc{JekelSandhuShi2024-PotentialHessianAscentSK}{Thm~4.1(3), pp.~30--31}
and \sloc{DaviesEtAl2026-PotentialHessianAscentIIISampling}{Cor.~B.35, p.~179}
against \ploc{Lemma~3.1 (p.~12)}{matrices.tex:114--134}.
\item \sloc{DaviesEtAl2026-PotentialHessianAscentIIISampling}{p.~17;
Lemma~B.18(b), p.~164} with $D=A^{-1}$, $\beta^2=j$, against
\ploc{$\mathsf P(K_A)=I$ and (3.18) (p.~16)}{matrices.tex:433--459}.
\item \sloc{JekelSandhuShi2024-PotentialHessianAscentSK}{Prop.~4.20,
Lemma~4.24, pp.~43--47} against \ploc{(3.15)--(3.18) (pp.~15--16)}{matrices.tex:356--436}.
\end{enumerate}

\space%
\finding[Planted Kac--Rice first moment for TAP zeros]{A planted Kac--Rice first moment for TAP zeros: determinant rate, Gaussian residual density, and the sign of the complexity functional}{planted-tap-kac-rice}

\paragraph{Summary.}
Fan, Mei and Montanari bound the expected number of TAP critical points of
the planted Sherrington--Kirkpatrick model by a Kac--Rice integral, evaluate
its rate as a squared discrepancy minus the relative entropy of the empirical
field law to a Gaussian, and show that on the Nishimori line this rate
vanishes at the Bayes-consistent law and, for large signal
($\lambda\ge\lambda_0$), is negative near it
\sloc{FanMeiMontanari2018-TAPFreeEnergyVariationalInference}{\S3, pp.~8--12; Prop.~5.2, p.~22}.
The preprint does not cite this paper and never uses the words Kac--Rice or
complexity: its introduction describes the argument only through its own
ingredients, ``A scalar entropy inequality and a local-volume estimate''
(\ploc{\S1.1, p.~4}{introduction.tex:122--125}), and calls the integral ``a
Gaussian integral'' (\ploc{\S1.3, p.~5}{introduction.tex:191--192}).
The computation is the smoothed integral \ploc{(5.5), p.~28}{stability.tex:119--134}
with \ploc{Lemma~5.3 (pp.~28--29)}{stability.tex:136--203}, the determinant
rate of \ploc{Lemma~3.5 (p.~20)}{matrices.tex:721--742}, and
\ploc{Lemma~4.1 (p.~22)}{entropy.tex:16--204} applied in
\ploc{(5.9), p.~29}{stability.tex:205--224}.
Both papers exclude TAP points with the wrong empirical field law by a first moment
over the field vector under the planted law, whose integrand is the Gaussian
density of the TAP residual times the determinant at its exponential rate and
whose rate in $y=\operatorname{arctanh}m$ vanishes at the Nishimori-consistent
Gaussian; the insight that this planted functional localizes the TAP points is
Fan, Mei and Montanari's.
\rel{adapted}: at $t=\sigma=0$ the preprint's rate functional is theirs
exactly, but the preprint proves a global inequality with its equality case
where they prove local negativity at large signal, smooths the residual, and
uses the determinant rate as a weight backed by a high-probability bound.
Celentano, Fan and Mei
\sloc{CelentanoFanMei2021-LocalConvexityTAPZ2Sync}{Lemma~4.7, p.~17; App.~B.2, pp.~41--43}
use one step of the approach, and Belius, Concetti and Genovese
\sloc{BeliusConcettiGenovese2024-DeterminantBrayMooreTAP}{Thms~1.1--1.2, pp.~2--3}
and Boursier \sloc{Boursier2026-LegendreStructureTAPComplexity}{Lemma~2.7, p.~20}
reprove the determinant rate; none of the three is cited.

\paragraph{Setup.}
Here $j=\beta^2<1$, and $F_{J,h_0}$ and $V_y$ are as in
Section~\ref{sec:overview}. The planted law is
$\tilde J=W+\frac jn\one\one^{\mathsf T}$, with $W$ a centered GOE of scale $j$,
and $h_0=t\one+\sqrt t\,g$ \ploc{(5.4), p.~28}{stability.tex:107--118}. For
$y\in\R^n$ put $m=\tanh y$, $q=n^{-1}\sum_im_i^2$, $b=1-q$, $B=jb$,
$M=n^{-1}\sum_im_i$ and, for $\sigma>0$,
\[
 d=t+jM,\quad s=t+\sigma^2+jq,\quad S=s+jq,\quad
 a_y=\frac1n\sum_i(y_i-d)m_i,\quad \psi=\frac{j(a_y-sb)^2}{2sS}
\]
\ploc{(5.6), p.~28}{stability.tex:136--145}; $N(d,s)$ has variance $s$,
$D(\cdot\Vert\cdot)$ is relative entropy, and $P$ is the empirical law of the
$y_i$ (or any law with finite second moment). The integral of
\ploc{(5.5)}{stability.tex:119--127} is
\[
 I_{t,\sigma}(\mathcal E)=(2\pi\sigma^2)^{-n/2}\int\one_{\{(y,\tilde J)\in\mathcal E\}}\,
 e^{-\norm{F_{\tilde J,h_0}(y)}^2/(2\sigma^2)+njb(y)^2/2}\,dy,
\]
and $F'=F'_{\tilde J,h_0}(y)=I+(BI-\tilde J)V_y-\frac{2j}nmm^{\mathsf T}V_y$
(\ploc{Lemma~5.6, p.~35}{stability.tex:719--720}; cf.\
\ploc{p.~38}{stability.tex:938--945}).
Fan, Mei and Montanari observe $Y=\frac\lambda nxx^{\mathsf T}+W'$ with
$W'\sim\mathrm{GOE}(n)$ (off-diagonal variance $1/n$) and take $x=\one$;
$g_n(m)$ and $H_n(m)$ are the scaled gradient and Hessian of the TAP free
energy at inverse temperature $\beta$ ((3.2)--(3.3), p.~8),
$Q(m)=\norm m^2/n$, $M(m)=\ip{\one}{m}/n$,
$A(m)=n^{-1}\sum_im_i\operatorname{arctanh}m_i$, $E(m)$ is a fourth statistic,
$\rho$ is the empirical law of the $m_i$, and $H(\cdot\,|\,\cdot)$ is relative
entropy. With $\beta=\lambda=\sqrt j$
(their Nishimori line), $\sqrt j\,W'=W$, $y=\operatorname{arctanh}m$ and $t=0$,
one has $g_n(m)=F_{\tilde J,0}(y)$ and $F'=H_n(m)V_y$.

\paragraph{Role in the argument.}
The integral bounds the expected weight of unstable configurations, and the
determinant rate turns an unstable approximate zero into a lower bound for
it: \ploc{Lemma~4.1}{entropy.tex:16--31} $\to$
\ploc{(5.9)}{stability.tex:205--214} $\to$
\ploc{Lemma~5.4}{stability.tex:226--389} $\to$
\ploc{(5.25), p.~37}{stability.tex:891--906}; with
\ploc{Lemma~3.5}{matrices.tex:721--742} and \ploc{Lemma~5.6}{stability.tex:713--832},
Markov's inequality excludes every witness
\ploc{(5.26)}{stability.tex:912--969} $\to$
\ploc{Proposition~5.1}{stability.tex:35--60} $\to$
\ploc{Corollary~5.2}{stability.tex:62--72} $\to$
\ploc{the preprint's \S2}{observation.tex} $\to$
\ploc{Theorem~1.1}{introduction.tex:33--42}. The neighbouring steps are
findings~\ref{f:conditional-hessian-first-moment} (the conditional matrix
factor in \ploc{Lemma~5.3}{stability.tex:157--171}),
\ref{f:approximate-zero-local-volume} (local volume),
\ref{f:nishimori-consistent-gaussian} (the equality law,
\ploc{Lemma~4.2}{entropy.tex:220--243}) and \ref{f:tap-stability-route} (the
route of Proposition~5.1).

\paragraph{The earlier technique.}
Fan, Mei and Montanari's Theorem~1.1 bounds the growth rate of
$\E[\mathrm{Crit}_n(U,V_n)]$, the expected number of TAP critical points with
$(Q,M,A,E)\in U\subseteq[\eta,1]\times\R^3$ and $m\in V_n^n$ ($V_n$ an interval
exhausting $(-1,1)$), by $\sup_US_\star$
\sloc{FanMeiMontanari2018-TAPFreeEnergyVariationalInference}{Thm~1.1, p.~4};
they present the proof as the Kac--Rice complexity strategy of earlier work
on spherical models (p.~4). In their \S3, the Kac--Rice bound
$\E[\mathrm{Crit}_n(U,V)]\le\int_{V^n}\one_U\,\E[|\det H_n(m)|\mid g_n(m)=0]\,p_m(0)\,dm$
(Prop.~3.1, p.~9), the determinant rate
$L(m)=\frac{\beta^2(1-Q(m))^2}{2}+\frac1n\sum_i\log\frac1{1-m_i^2}$
(Prop.~3.2, p.~9, proved in their \S4 through the free convolution) and the
gradient density, a Gaussian mixture over one auxiliary scalar (their $y$)
(Lemma~3.3, p.~9), give the integral $T(U,V_n)$ of Cor.~3.4 (pp.~9--10).
Sanov's theorem and a Varadhan bound (Prop.~3.5, p.~10), then the supremum
over their $y$, give the rate
$\frac1{4\beta^2}(\beta^2(1-q)+\beta\lambda\varphi^2/q-a/q)^2-H(\rho\,|\,\rho_{\varphi,q})$
(\S3.4, p.~11), with $\rho_{\varphi,q}$ the law of $\tanh z$, $z\sim N(\beta\lambda\varphi,\beta^2q)$;
Lagrange multipliers for the constraints give $S_\star$ (p.~12). On the
Nishimori line $\beta=\lambda$, Proposition~5.2 (p.~22) shows
$S_\star(q_\star,\varphi_\star,a_\star,e_\star)=0$ for every $\lambda>0$ and,
for $\lambda\ge\lambda_0$ and $|q-q_\star|,|\varphi-\varphi_\star|\le c_0/\lambda^2$,
$S_\star\le-c(\lambda)\norm{(q,\varphi,a,e)-(q_\star,\varphi_\star,a_\star,e_\star)}^2$;
here $q_\star=0$ when $\lambda\le1$ (Prop.~5.5, p.~23).

\paragraph{Correspondence.}
\begin{center}\small
\begin{xltabular}{\linewidth}{@{}>{\raggedright\arraybackslash}p{0.15\linewidth}>{\raggedright\arraybackslash}X>{\raggedright\arraybackslash}X@{}}
\toprule
 & Preprint & Fan--Mei--Montanari \cite{FanMeiMontanari2018-TAPFreeEnergyVariationalInference}\\
\midrule
\endfirsthead
\toprule
 & Preprint & Fan--Mei--Montanari \cite{FanMeiMontanari2018-TAPFreeEnergyVariationalInference}\\
\midrule
\endhead
Determinant rate & $\frac1n\log\det(F'^{\mathsf T}F'+\gamma I)\le jb^2+\varepsilon$ with high probability, simultaneously over $y$ \ploc{(Lemma~3.5, (3.26))}{matrices.tex:721--742} & $\E[|\det H_n(m)|\mid g_n(m)=0]\le e^{nL(m)+o(n)}$ for $Q(m)\ge\eta$, $m\in V_n^n$ (Prop.~3.2, (3.5), p.~9)\\
\addlinespace
First moment & $\E I_{t,\sigma}(\mathcal E)$, weight $e^{njb^2/2}$ \ploc{(5.5)}{stability.tex:119--134} & $T(U,V_n)=e^{o(n)}\int_{(-1,1)^n}\one_U\,e^{nL(m)}p_m(0)\,dm$ (Cor.~3.4, pp.~9--10)\\
\addlinespace
Residual density & $\Cov(Wm)=jqI+\frac jnmm^{\mathsf T}$; $\E I_{t,\sigma}=\int\sqrt{s/S}\prod_i\varphi_{d,s}(y_i)\,e^{n\psi}\cdots dy$ \ploc{(Lemma~5.3)}{stability.tex:146--203} & $\E[(W'm)_i(W'm)_k]=Q\delta_{ik}+\frac{m_im_k}n$ (Lemma~3.3, p.~9); supremum over their $y$ (\S3.4, p.~11)\\
\addlinespace
Rate & $\psi-D(P\Vert N(d,s))$ \ploc{(5.6), (5.9)}{stability.tex:136--214} & squared term $-\,H(\rho\,|\,\rho_{\varphi,q})$ (\S3.4, p.~11)\\
\addlinespace
Sign & $D(P\Vert N(d,s))\ge\psi$ for every $P$ with finite second moment when $s\ge jq$; equality iff $P=N(d,s)$ \ploc{(Lemma~4.1, (4.1))}{entropy.tex:16--31} & $S_\star=0$ at the Bayes point; local negativity for $\lambda\ge\lambda_0$ (Prop.~5.2, (5.4), p.~22)\\
\addlinespace
\bottomrule
\end{xltabular}
\end{center}
At $t=\sigma=0$, $\beta=\lambda=\sqrt j$ and $\varphi=M$, one has $d=jM=\beta\lambda\varphi$,
$s=jq=\beta^2q$, $S=2jq$ and $a_y=a-jM^2$ with $a=n^{-1}\sum_iy_im_i=A(m)$, so
\[
 \psi=\frac{(a-jM^2-jq(1-q))^2}{4jq^2}
 =\frac1{4\beta^2}\Bigl(\beta^2(1-q)+\frac{\beta\lambda\varphi^2}{q}-\frac aq\Bigr)^2,
 \qquad
 D(P\Vert N(jM,jq))=H(\rho\,|\,\rho_{\varphi,q}),
\]
the second because $\tanh$ is a bijection carrying $N(\beta\lambda\varphi,\beta^2q)$
to $\rho_{\varphi,q}$. For the determinant, $F'=H_n(m)V_y$ and
$dm=\prod_i(1-m_i^2)\,dy$ give $e^{nL(m)}\,dm=e^{njb^2/2}\,dy$: the rate
$jb^2/2$ of $\frac1n\log|\det F'|$ in \ploc{Lemma~3.5}{matrices.tex:721--742}
is $L(m)$ without its Jacobian term $\frac1n\sum_i\log\frac1{1-m_i^2}$, and
$T(U,V_n)$ is, up to $e^{o(n)}$, the integral
$\int_{\R^n}\one_U\,e^{njb^2/2}p_{\tanh y}(0)\,dy$, the $\sigma\to0$ limit of
$\E I_{0,\sigma}$ restricted to $(Q,M,A,E)\in U$.

\paragraph{Differences.}
\begin{itemize}
\item The sign is proved for every $P$ with finite second moment whenever
$s\ge jq$, with its equality case, by the transport $T_u(y)=y-u\tanh y$
\ploc{(4.4), p.~23}{entropy.tex:62--83}; Fan, Mei and Montanari's negativity
holds only within $c_0/\lambda^2$ of the Bayes point, for $\lambda\ge\lambda_0$
and $Q\ge\eta>0$.
\item The regime is $j<1$ (their $\beta=\lambda<1$, where $q_\star=0$), with
$t\ge0$ and $\sigma>0$; the Gaussian reduction is an exact identity with
prefactor $\sqrt{s/S}$, not a supremum over their $y$.
\item \ploc{Lemma~3.5}{matrices.tex:721--742} is an unconditional
high-probability bound on $\log\det(F'^{\mathsf T}F'+\gamma I)$ for the
non-symmetric $F'$, simultaneous over diagonal $V$ and finite-rank
perturbations and proved by interpolation in $z$
\ploc{(3.27)--(3.29)}{matrices.tex:743--812}; Fan, Mei and Montanari bound
$\E[|\det H_n(m)|\mid g_n(m)=0]$ at each fixed $m$.
\end{itemize}

\paragraph{How to check.}
\begin{enumerate}
\item \ploc{(5.6)}{stability.tex:136--145} and \ploc{(5.9)}{stability.tex:205--214}
at $t=\sigma=0$ against
\sloc{FanMeiMontanari2018-TAPFreeEnergyVariationalInference}{\S3.4, p.~11},
via the display after the table.
\item The proof of \ploc{Lemma~5.3}{stability.tex:173--203} against
\sloc{FanMeiMontanari2018-TAPFreeEnergyVariationalInference}{Lemma~3.3 and Cor.~3.4, pp.~9--10}
after the substitution $m=\tanh y$.
\item \ploc{Lemma~3.5}{matrices.tex:721--742}, with $F'$ as on
\ploc{p.~38}{stability.tex:938--947}, against
\sloc{FanMeiMontanari2018-TAPFreeEnergyVariationalInference}{Prop.~3.2, (3.5), p.~9}
using $F'=H_n(m)V_y$.
\item \ploc{Lemma~4.1 (p.~22)}{entropy.tex:16--31} against
\sloc{FanMeiMontanari2018-TAPFreeEnergyVariationalInference}{Prop.~5.2, p.~22}.
\end{enumerate}

\space%
\finding[Approximate zeros from a smoothed first moment]{Excluding approximate zeros by a smoothed first moment and local volume}{approximate-zero-local-volume}

\paragraph{Summary.}
Huang and Sellke
\sloc{HuangSellke2023-StrongTopologicalTrivialization}{\S5, Thm~5.2, Lemma~5.11, pp.~28--36}
(2023) show how an annealed bound on exact critical points with
a bad property rules out, with probability $1-e^{-cN}$, approximate ones with
that property; the preprint cites them nowhere, calls the step a
``local-volume estimate'' (\ploc{\S1.1, p.~4}{introduction.tex:122--125}) and
a ``local-volume lower bound'' (\ploc{\S1.3, p.~5}{introduction.tex:191--194})
without citation, and never mentions Kac--Rice.
Their insight: approximate zeros cannot be counted, so one integrates a narrow
Gaussian in the residual, weighted by the Jacobian determinant (in Huang and
Sellke, its conditional expectation, of known exponential rate), over the bad
set and bounds its mean; near an approximate zero the Gaussian volume from the
linearization is $|\det|^{-1}$, cancelling the weight wherever the determinant
is at most its typical rate (everywhere at once, with high probability), so
the integral is at least $e^{-o(n)}$, and Markov's inequality concludes.
The preprint follows this in \ploc{(5.5)}{stability.tex:119--134},
\ploc{Lemma~5.6}{stability.tex:713--832} and the proof of
\ploc{Proposition~5.1}{stability.tex:908--969}, with
\ploc{Lemma~3.5}{matrices.tex:721--869} as the uniform determinant bound;
\rel{adapted}, because only the external field is rerandomized, the
deterministic weight $e^{njb(y)^2/2}$ replaces the determinant, and
regularizing the Jacobian $L$ to $L^{\mathsf T}L+\gamma I$ replaces their
spectral split.
Huang, Montanari and Pham
\sloc{HuangMontanariPham2024-SphericalSpinGlassSamplingSL}{p.~29}, also not
cited, use a consequence of Theorem~5.2 (finding~\ref{f:tap-stability-route}).

\paragraph{Correspondence.}
\begin{center}\small
\begin{xltabular}{\linewidth}{@{}>{\raggedright\arraybackslash}p{0.12\linewidth}>{\raggedright\arraybackslash}X>{\raggedright\arraybackslash}X@{}}
\toprule
 & Preprint & Huang--Sellke\\
\midrule
\endfirsthead
\toprule
 & Preprint & Huang--Sellke\\
\midrule
\endhead
Smoothed integral & \ploc{(5.5), p.~28}{stability.tex:119--134}: $I_{t,\sigma}(\mathcal E)=(2\pi\sigma^2)^{-n/2}\int_{\mathcal E}e^{-\norm{F(y)}^2/(2\sigma^2)+njb(y)^2/2}\,dy$, $F=F_{\tilde J,h_0}$; by the area formula, the expected number of zeros of $F_{\tilde J,h_0+\sigma\xi}$ in $\mathcal E$, with $e^{njb(y)^2/2}$ for $|\det L|$ & given $H_N$, the density at $0$ of $\nabla_{\mathrm{sp}}H_{N,\delta}(y)$, $H_{N,\delta}=\sqrt{1-\delta}\,H_N+\sqrt\delta\,H_N'$, times $\E[|\det\nabla^2_{\mathrm{sp}}H_{N,\delta}(y)|\mid\cdots]=e^{N\int\log|\lambda|\,d\mu\pm o(N)}$ ((5.20)--(5.22), pp.~33--34)\\
\addlinespace
Local lower bound & \ploc{Lemma~5.6}{stability.tex:713--742}: $\norm{F(y)}\le\rho\sqrt n$ and $\frac1n\log\det Q\le jb(y)^2+\varepsilon_1$, $Q=L^{\mathsf T}L+\gamma I$, give $\ge e^{-\varepsilon n}$ over a ball \ploc{(5.19)--(5.20), p.~35}{stability.tex:723--734}; the ellipsoid $y+\sigma Q^{-1/2}g''$ gives $(\det Q)^{-1/2}$ \ploc{(5.21), p.~36}{stability.tex:745--759}, cancelling $e^{njb(y)^2/2}$ (\ploc{p.~37}{stability.tex:818--831}) & Lemma~5.11 (p.~33): from $\norm{\nabla_{\mathrm{sp}}H_N(x)}\le\varepsilon\sqrt N$ (5.15), $\ge e^{-o_\delta(N)}$ expected critical points of $H_{N,\delta}$ in $B_{\alpha\sqrt N}(x)$; the volume $\mathrm{Vol}(S_1)|\det(J_N(x)|_{S_2})|^{-1}$ ((5.23), p.~35) cancels (5.21)\\
\addlinespace
Uniform determinant & \ploc{Lemma~3.5}{matrices.tex:721--742}: with probability $\ge1-Ce^{-c_4n}$, $\frac1n\log\det(L^{\mathsf T}L+\gamma I)\le jb^2+\varepsilon$ for all diagonal $0\le V\le I$ and bounded rank-$r_0$ $E$ at once \ploc{(3.26), p.~20}{matrices.tex:721--742}; used at every $y$ (\ploc{p.~38}{stability.tex:938--957}) & $H_N\in K_N(\varepsilon)$ (Def.~13(c), Prop.~5.8, p.~32): with probability $\ge1-e^{-c(\varepsilon)N}$, for all $x$ the spectrum of $J_N(x)$ is $W_1$-close to $\tilde\mu_{\xi,\nabla_{\mathrm{rad}}H_N(x)}$\\
\addlinespace
Conclusion & a witness \ploc{(5.26), p.~37}{stability.tex:908--920} forces $I_{t,\sigma}(E_t)\ge e^{-\varepsilon n}$ for the unstable set $E_t$; Markov and $\E I_{t,\sigma}(E_t)\le e^{-a_0n}$ \ploc{(5.25), p.~37}{stability.tex:892--896} bound its probability by $Ce^{-a_2n}$ (\ploc{p.~38}{stability.tex:958--969}) & (5.24), averaged, and $\E|\mathsf{Crt}_N^{\overline{\mathcal D},0,\upsilon/2}(H_N)|\le e^{-c_0N}$ (5.3) bound the probability of a bad approximate critical point by $e^{-cN}$ (p.~36)\\
\bottomrule
\end{xltabular}
\end{center}

\paragraph{How to check.}
\begin{enumerate}
\item \ploc{(5.5) and Lemma~5.6, pp.~28 and 35--37}{stability.tex:119--134, 713--832}
against \sloc{HuangSellke2023-StrongTopologicalTrivialization}{(5.20)--(5.23), Lemma~5.11, pp.~33--35}.
\item \ploc{Lemma~3.5}{matrices.tex:721--869} and
\ploc{pp.~37--38}{stability.tex:908--969} against
\sloc{HuangSellke2023-StrongTopologicalTrivialization}{Def.~13(c), Prop.~5.8, p.~32; (5.24), p.~36}.
\end{enumerate}

\space%
\finding{TAP stability along a localization path by a planted first moment}{tap-stability-route}

\paragraph{Summary.}
Huang, Montanari and Pham prove, for spherical spin glasses at a fixed time
of the stochastic-localization path, that with probability $1-e^{-cN}$ the
TAP free energy has a unique well-conditioned critical point in a band, near
which every approximate critical point lies
\sloc{HuangMontanariPham2024-SphericalSpinGlassSamplingSL}{Prop.~6.1, pp.~26--27};
the preprint does not cite that paper.
Their route gets well-posedness from a first moment under the planted law of
El Alaoui, Montanari and Sellke: bound the bad critical points by a
conditional probability of a bad Hessian times a subexponential total, pass
to approximate critical points, and transfer back through the normalized
partition function.
The preprint follows this route for the field equation at the fields its
observation path visits, in \ploc{Proposition~5.1 (p.~27)}{stability.tex:35--60}
with \ploc{Lemma~5.7 (p.~40)}{stability.tex:1105--1125}; its introduction
attributes this uniformity to its scalar entropy inequality and local-volume
estimate (\ploc{\S1.1, p.~4}{introduction.tex:122--125}) without naming an
earlier instance of the route.
The relationship is \rel{adapted}: the total comes from the entropy
inequality, not topological trivialization; the bound is uniform in $t\le T$
through a fixed mesh; uniqueness of the root is global and deterministic;
the transfer is quenched, not by contiguity; and approximate roots
are handled neither by the theorem of Huang and Sellke that Huang, Montanari
and Pham cite nor by the device behind it, which rerandomizes the disorder
\sloc{HuangSellke2023-StrongTopologicalTrivialization}{(5.1) and Thm~5.2, p.~28},
but by the preprint's own device on the same principle (an approximate zero
forces annealed mass $e^{-o(n)}$): a residual-smoothed integral, in effect a
perturbation of the field, with a local-volume lower bound
(finding~\ref{f:approximate-zero-local-volume}).

\paragraph{Correspondence.}
\begin{center}\small
\begin{xltabular}{\linewidth}{@{}>{\raggedright\arraybackslash}p{0.13\linewidth}>{\raggedright\arraybackslash}X>{\raggedright\arraybackslash}X@{}}
\toprule
Step & Preprint & Huang--Montanari--Pham\\
\midrule
\endfirsthead
\toprule
Step & Preprint & Huang--Montanari--Pham\\
\midrule
\endhead
Conclusion & with $F=F_{J,Y(t)}$, \ploc{(5.3), p.~27}{stability.tex:45--56}: for $J\in\mathcal J_n$, conditional probability $\ge1-e^{-an}$, all $t\le T$, $\norm{F(y)}\le2\rho_0\sqrt n\Rightarrow\mathsf H_J(y,A)\succ cI$ for diagonal $A$ near $V_y$; \ploc{(5.29), p.~40}{stability.tex:1117--1121}: one zero $r$, and $\norm{y-r}\le C\norm{F(y)}$ if $\norm{F(y)}<2\rho_0\sqrt n$ & Prop.~6.1, pp.~26--27: under $\mathbb P$, one $t$, probability $1-e^{-cN}$, a unique critical point $m^{\mathrm{TAP}}$ in the band $\mathcal S_\iota$, Hessian spectrum in $[-C_{\max},-C_{\min}]$; every $\iota$-approximate critical point in $\mathcal S_\iota$ is within $o_\iota(1)$ of it in $\norm{\cdot}_N=\norm{\cdot}/\sqrt N$\\
\addlinespace
Planted first moment & \ploc{(5.25), p.~37}{stability.tex:891--906}, under \ploc{(5.4), p.~28}{stability.tex:107--118}: $\E I_{t,\sigma}(E_t)\le e^{-a_0n}$ for the smoothed integral over unstable $y$, a conditional probability $e^{-a_*n}$ (\ploc{Lemma~5.5}{stability.tex:403--423}) times a scalar integral $e^{o(n)}$ (\ploc{Lemma~5.4(i)}{stability.tex:229--238}) & under $\mathbb P$ ((3.2), p.~12), Kac--Rice (6.6) and Cauchy--Schwarz: $\E|\mathrm{Crt}\cap\mathcal T|\le C\sup\P[m\in\mathcal T\mid\cdots]^{1/2}$, using $\E|\mathrm{Crt}|=2+o_N(1)$ (Lemma~6.8, Prop.~6.7, p.~30; Prop.~6.9, p.~31)\\
\addlinespace
Approximate roots & a witness \ploc{(5.26)}{stability.tex:911--916} forces $I_{t,\sigma}(E_t)\ge e^{-\varepsilon n}$ (\ploc{Lemma~5.6}{stability.tex:713--832}); Markov's inequality (\ploc{p.~38}{stability.tex:938--969}) & Huang and Sellke's Theorem~1.6 \sloc{HuangSellke2023-StrongTopologicalTrivialization}{p.~8}, in the proof of Prop.~6.2(a) (p.~29)\\
\addlinespace
Transfer & $L(J)=Z(J)/\E Z(J)$, lower tail \ploc{(5.27), p.~39}{stability.tex:1040--1065}, Markov's inequality (\ploc{pp.~39--40}{stability.tex:1067--1083}) & $d\mathbb P/d\mathbb Q=Z(H_N)/\E Z(H_N)$, contiguity (Prop.~3.4, Cor.~3.5, p.~12; used p.~93)\\
\bottomrule
\end{xltabular}
\end{center}

\paragraph{How to check.}
\begin{enumerate}
\item \ploc{(5.3), p.~27}{stability.tex:45--56} and
\ploc{(5.29), p.~40}{stability.tex:1117--1121} against
\sloc{HuangMontanariPham2024-SphericalSpinGlassSamplingSL}{Prop.~6.1, pp.~26--27; (4.8)--(4.9), p.~15}.
\item The plan on \ploc{p.~27}{stability.tex:99--105} and
\ploc{\S5.4}{stability.tex:836--1094} against \S3.2 and the proofs of
Props.~6.1--6.2 \sloc{HuangMontanariPham2024-SphericalSpinGlassSamplingSL}{pp.~11--12, 27--32};
for approximate roots (\ploc{pp.~37--38}{stability.tex:908--969}), also against
\sloc{HuangSellke2023-StrongTopologicalTrivialization}{(5.1), Thm~5.2, p.~28}.
\end{enumerate}

\space%
\finding[Conditional TAP Hessian near the Bayes law]{Conditioning the TAP Hessian on the residual near the Bayes law to exclude unstable zeros}{conditional-hessian-first-moment}

\paragraph{Summary.}
Celentano, Fan and Mei prove local strong convexity of the TAP free energy
for $\mathbb Z_2$ synchronization by ``a second Kac--Rice upper bound''
\sloc{CelentanoFanMei2021-LocalConvexityTAPZ2Sync}{\S4.1, p.~14}: the event
that the Hessian's least eigenvalue near $m$ is small enters the count of
critical points as a probability conditional on a vanishing gradient
(Lemma~4.8); conditioning leaves a fresh projected GOE plus explicit
low-rank terms (Lemma~B.2); and near the Bayes (Nishimori) law this
probability is exponentially small, by a finite-dimensional Schur complement
whose positivity rests on $1-\lambda^2\E[(1-m^2)^2]>0$ (Lemma~4.9,
Cor.~4.10, p.~18). The preprint does not cite this paper, and cites no
other work for this step. The marked Kac--Rice bound itself is generic: they
derive it from Fan, Mei and Montanari's arguments and a theorem of Adler and
Taylor \sloc{CelentanoFanMei2021-LocalConvexityTAPZ2Sync}{App.~A.5,
pp.~33--34}. What is theirs is its combination with the TAP equation and the
Bayes law. The preprint follows that combination: in an exact Gaussian
identity for a smoothed count of approximate zeros, the instability event
enters as a probability conditional on the residual
\ploc{(5.8), p.~28}{stability.tex:157--203}, and
\ploc{Lemma~5.5 (pp.~31--35)}{stability.tex:393--700} bounds it by
$e^{-a_*n}$ near the equality laws $N(s_*,s_*)$
(finding~\ref{f:nishimori-consistent-gaussian}), giving
\ploc{(5.25), p.~37}{stability.tex:864--906}. \rel{adapted}: the
conditioning is on a noisy residual, their Sudakov--Fernique variational
bound becomes whitening plus an exact Schur complement whose determinant
equals theirs up to a positive factor, and positivity is simultaneous over
nearby diagonal weights; Huang, Montanari and Pham later take the same step
for spherical TAP along stochastic localization
\sloc{HuangMontanariPham2024-SphericalSpinGlassSamplingSL}{Lemma~6.8,
Prop.~6.9, pp.~30--31}.

\paragraph{Correspondence.}
\begin{center}\small
\begin{xltabular}{\linewidth}{@{}>{\raggedright\arraybackslash}p{0.14\linewidth}>{\raggedright\arraybackslash}X>{\raggedright\arraybackslash}X@{}}
\toprule
 & Preprint & Celentano--Fan--Mei\\
\midrule
\endfirsthead
\toprule
 & Preprint & Celentano--Fan--Mei\\
\midrule
\endhead
First moment & $\E I_{t,\sigma}(\mathcal E)=\int\sqrt{s/S}\,G_{t,\sigma,n}(y)\,\P\bigl((y,\tilde J)\in\mathcal E\mid Wm+\sqrt{t+\sigma^2}\,g'=z\bigr)\,dy$ \ploc{(5.8), p.~28}{stability.tex:157--166} & $\E|\mathcal C|\le\int_T\E\bigl[|\det H(m)|\,\one\{\ell(m,W)+U<t\}\bigm|g(m)=0\bigr]\,p_{g(m)}(0)\,dm$ ((4.18), p.~18)\\
\addlinespace
Conditional law & $W=P_uW_0P_u+uw^{\mathsf T}+wu^{\mathsf T}+\alpha uu^{\mathsf T}$, $u=m/\norm m$, $w=\bar w+\kappa P_uW_0u$, $W_0$ a fresh GOE \ploc{(5.12), p.~32}{stability.tex:454--463} & given $g(m)=0$, $W$ is $P_m^\perp\tilde WP_m^\perp$, $\tilde W$ a fresh GOE, plus rank-two terms in $m$, $z(m)$ (Lemma~B.2, pp.~43--44)\\
\addlinespace
Bayes law & $(P_y,t)$ near $\{(P_t,t)\}$, $P_y$ the empirical law of $y$, $P_t=N(s_*,s_*)$ (\ploc{Lemma~5.5, p.~31}{stability.tex:399--423}) & $\mathcal D_\eta$: $W(\hat\mu_m,\mu_\star)<\eta$, with $\operatorname{arctanh}m\sim N(\lambda^2q_\star,\lambda^2q_\star)$ under $\mu_\star$ ((4.5)--(4.8), p.~14)\\
\addlinespace
Sign & Gram matrix $H$ \ploc{(5.16), p.~33}{stability.tex:549--573}; $\det(I-HC)=(d_0+2g_1)(d_0+g_1)^2>0$ \ploc{(5.17), p.~34}{stability.tex:579--587} & Schur complement along $m$, $\one$: $c_2>0$, $c_1-c_2>0$ ((B.13)--(B.14), pp.~39--40, 48--49)\\
\bottomrule
\end{xltabular}
\end{center}

\paragraph{How to check.}
\begin{enumerate}
\item \ploc{(5.8), p.~28}{stability.tex:157--166} and the end of its proof
(\ploc{p.~29}{stability.tex:196--202}) against
\sloc{CelentanoFanMei2021-LocalConvexityTAPZ2Sync}{Lemma~4.8, pp.~17--18}.
\item \ploc{(5.12), p.~32}{stability.tex:454--463} at $t=\sigma=0$ against
\sloc{CelentanoFanMei2021-LocalConvexityTAPZ2Sync}{Lemma~B.2, pp.~43--44},
with $\lambda^2=j$, $\lambda W\leftrightarrow W$, $z=z(m)$.
\item \ploc{(5.16)--(5.17), pp.~33--34}{stability.tex:549--587} against
\sloc{CelentanoFanMei2021-LocalConvexityTAPZ2Sync}{(B.13)--(B.14),
pp.~39--40}: with $\lambda^2=j$,
$c_2(c_1-c_2)=(1-jf)^2(1-jf+je)/(fe)=\det(I-HC)/(fe)$, where
$f=\E[(1-m^2)^2]$ and $e=\E[m^2(1-m^2)]$ under the Bayes law.
\end{enumerate}

\space%
\finding[Memory corrections for every input, cancelling one-step reversals]{A memory correction for every input of a multi-input recursion, and its cancellation of one-step reversals at every depth}{amp-memory-onsager-rule}

\paragraph{Summary.}
Bayati, Lelarge and Montanari (BLM) correct each input of a multi-input
nonlinearity in approximate message passing (AMP) by its averaged partial
derivative times that input's own source
\sloc{BayatiLelargeMontanari2012-UniversalityMessagePassing}{(1.4)--(1.6),
p.~7}, and prove that these memory terms cancel, at every depth, the one-step
reversing contributions to the expansion of the iterates over labelled trees
(\S2, \S4); Montanari \sloc{Montanari2018-OptimizationSKHamiltonian}{(2.1),
p.~6} writes the rule for every earlier iterate. The preprint cites neither
paper: it cites Thouless, Anderson and Palmer for the Onsager correction
\ploc{(\S1.1, p.~4)}{introduction.tex:125--128}, credits (6.1) to Bolthausen
\ploc{(\S6.1, p.~42)}{residuals.tex:32--34}, whose recursion has a single
memory term with a deterministic coefficient
\sloc{Bolthausen2012-IterativeTAPConstruction}{p.~4}, and names AMP only
in \S1.1, citing El Alaoui, Montanari and Sellke and Celentano for other
content. Its rule (6.7) in
\ploc{Definition~6.2 (p.~44)}{residuals.tex:179--222} is Montanari's (2.1)
term for term after rescaling; it is used again for the implicit field
\ploc{(6.14), p.~46}{residuals.tex:343--362}, and in the proof of
\ploc{Lemma~6.5, (6.33), p.~54}{residuals.tex:992--1003}, and the
cancellation runs through the proofs of
\ploc{Lemmas~6.1 (p.~43)}{residuals.tex:105--127},
\ploc{6.4 (pp.~51--52)}{residuals.tex:750--842} and
\ploc{3.3 (p.~16)}{matrices.tex:438--459}. The shared insight, made rigorous
by BLM: in a recursion in a symmetric random matrix, the averaged derivative
times an input's source is minus the leading contribution of a step that
immediately reverses, so subtracting it for each input frees each coordinate
from its own back-reaction. The relationship is \rel{adapted}: the preprint
cancels the diagonal noncrossing-pairing predictions of formal spin-flip
Jacobians, not moments of the iterates, by a signed
inclusion--exclusion over adjacent pairs, not by comparison with
non-backtracking iterates; Montanari's paper holds the rule alone.

\paragraph{Correspondence.}
\begin{center}\small
\begin{xltabular}{\linewidth}{@{}>{\raggedright\arraybackslash}p{0.1\linewidth}>{\raggedright\arraybackslash}X>{\raggedright\arraybackslash}X@{}}
\toprule
 & Preprint & BLM; Montanari\\
\midrule
\endfirsthead
\toprule
 & Preprint & BLM; Montanari\\
\midrule
\endhead
Rule & $y^a=Jw^a-j\sum_l\langle\partial_lw^a\rangle m^{l-1}$ $-j\sum_{b<a}\langle\partial_{y^b}w^a\rangle w^b$ \ploc{(6.7), p.~44}{residuals.tex:200--205}, $\langle\cdot\rangle$ the coordinate average & $u^{k+1}=Af_k(u^0,\dots,u^k)-\sum_{r=1}^kb_{k,r}f_{r-1}$, $b_{k,r}=n^{-1}\sum_i\partial f_k/\partial u^r_i$ (Montanari (2.1)); $-\sum_kA_{ik}^2\,\partial_xf^k(x^t_k)\,f^i(x^{t-1}_i)$ (BLM (1.6))\\
\addlinespace
Linear & formal Jacobians (6.3) with every diagonal $I$: $H_{l+1}=JH_l-jH_{l-1}$, $H_2=J^2-jI$ \ploc{(p.~42)}{residuals.tex:79--84} & $f(x)=x$: $x^{t+1}=Ax^t-x^{t-1}$, $x^2=A^2\one-\one$, cancelling the paths $1\to i\to1$ (p.~11)\\
\addlinespace
One step & $-jb_sI$ is minus the contraction of $JD_{1-(m^s)^2}J$ \ploc{(p.~43)}{residuals.tex:110--112}; an example \ploc{on p.~51}{residuals.tex:753--757} & $f(x)=x^3$: $b_1$ cancels the trees with a label equal to the root's ((2.6), pp.~11--12)\\
\addlinespace
Every depth & each subset $T$ of the chain's adjacent pairs contracted, sign $(-1)^{|T|}$: per full noncrossing pairing, $(1-1)^a=0$ \ploc{(pp.~43, 52)}{residuals.tex:117--124, 814--840} & $x^t_i=z^t_i$ (non-backtracking trees, Lemma~1, p.~19) plus trees with a backtracking 3-path or star (Lemma~3, p.~26; (4.29), p.~27), negligible by (4.30), p.~28\\
\addlinespace
\bottomrule
\end{xltabular}
\end{center}

\paragraph{How to check.}
\begin{enumerate}
\item \ploc{(6.7), p.~44}{residuals.tex:200--205} against
\sloc{Montanari2018-OptimizationSKHamiltonian}{(2.1), p.~6} with
$J=\sqrt j\,A$, fields divided by $\sqrt j$; and
\sloc{BayatiLelargeMontanari2012-UniversalityMessagePassing}{(1.5)--(1.6), p.~7}.
\item The proof of \ploc{Lemma~6.1 (p.~43)}{residuals.tex:105--127} with all
diagonals $I$ against
\sloc{BayatiLelargeMontanari2012-UniversalityMessagePassing}{(2.3), pp.~10--11}.
\item ``Cancellation of the formal traces'' \ploc{(pp.~51--52)}{residuals.tex:750--842}
against \sloc{BayatiLelargeMontanari2012-UniversalityMessagePassing}{Lemma~3,
(4.29)--(4.30), pp.~26--29}.
\end{enumerate}

\space%
\finding[Terminal gap via small principal submatrices]{A gap at a large observation time from few weak-field coordinates and uniformly small principal submatrices}{terminal-gap-sparse-submatrices}

\paragraph{Summary.}
Davies, Lee, Sandhu and Shi's Potential Hessian Ascent~III (PHA~III, May
2026) studies the SK posterior $\mu_{\beta A,y_T}$ at a large constant
localization time, the preprint's $\mu_{Y(T)}$
($A\sim\mathrm{GOE}(n)$, $J=\beta(A-\operatorname{diag}A)$ in law), and
states the insight \sloc{DaviesEtAl2026-PotentialHessianAscentIIISampling}{p.~19}:
only $O(\varepsilon n)$ coordinates can still fluctuate, for any prior
(Lemma~7.43), so only small principal submatrices of the disorder matter, and
a Gaussian norm tail with a union bound over supports makes them all small at
once (Lemma~7.23). The preprint cites neither PHA~III nor Davies, Lee, Sandhu
and Shi's weak Poincar\'e paper, whose Theorem~5.1
\sloc{DaviesEtAl2026-WeakPoincareApproxSL}{p.~21} uses the same two lemmas
for a weak Poincar\'e inequality for Glauber dynamics at
$\mu_{\beta A,y_T}$, $\beta<1/2$; at the lead-in to
\ploc{Lemma~8.2 (p.~61)}{terminal.tex:217--223} it cites only Bandeira, El
Alaoui and R\"odder's later Theorem~1.1
\cite{BandeiraElAlaouiRodder2026-GlauberHighOverlapGibbs}, for combining
small-submatrix bounds with a large constant field. Lemma~8.2 and the proof of
\ploc{Proposition~8.3 (pp.~63--64)}{terminal.tex:334--433} carry out the
mechanism at $Y(T)$ and obtain $\Var_{\mu_{Y(T)}}f\le2\D_{Y(T)}(f)$ for any
law of the spins (\rel{adapted}: both probabilistic steps also appear in
PHA~III, but its functional-inequality step on a Hamming ball, through the
needle decomposition of $\beta A+\gamma I$ that forces $\beta<1/2$, is
replaced by a split of the signed curvature form over a
configuration-dependent weak set).

\paragraph{Correspondence.}
\begin{center}\small
\begin{xltabular}{\linewidth}{@{}>{\raggedright\arraybackslash}p{0.15\linewidth}>{\raggedright\arraybackslash}X>{\raggedright\arraybackslash}X@{}}
\toprule
 & Preprint & PHA~III; weak Poincar\'e paper\\
\midrule
\endfirsthead
\toprule
 & Preprint & PHA~III; weak Poincar\'e paper\\
\midrule
\endhead
Why small submatrices & the few coordinates with appreciable conditional variance occupy a small principal submatrix \ploc{(p.~61)}{terminal.tex:217--218} & on a wedge only $O(\varepsilon n)$ coordinates can fluctuate, so small principal submatrices matter, not $A$ (p.~19)\\
\addlinespace
Few weak coordinates, any prior & $X$ of any law, $Y=Y(T)$: given $X$, the events $\{|Y_i|<T/2\}$ are independent with probability $p_T\le e^{-T/8}$; exponential Markov: at most $\alpha n/2$ occur, off probability $e^{-\alpha n/4}$ \ploc{(p.~63)}{terminal.tex:370--384} & $x_0\sim\mu$ of any law, $y_t=tx_0+B_t$: the events $\{(x_0)_i\ne\operatorname{sign}(y_t)_i\}$ are independent with probability $\Phi(-\sqrt t)$; Chernoff (Lemma~7.43, p.~135)\\
\addlinespace
All small submatrices at once & $\norm{U_{SS}},\norm{Q_{SS}}\le\varepsilon$ for all $|S|\le\alpha n$, where $U=\tanh J$, $Q=U\circ U$ \ploc{(8.11), p.~62}{terminal.tex:235--239}; for $U$, a net on each support and a union bound over the supports of size $\lfloor\alpha n\rfloor$ \ploc{(pp.~62--63)}{terminal.tex:297--314} & $\P(\forall|S|\le\varepsilon n:\norm{A_{S\times S}}\le c_1\sqrt{h(\varepsilon)})\ge1-2e^{-h(\varepsilon)n}$, $h$ the binary entropy, by the GOE norm tail and a union bound over supports (Lemma~7.23, p.~115)\\
\addlinespace
At the terminal measure & $S(x)=\{i:|Y_i+(Jx)_i|<T/4\}$ has $|S(x)|\le\alpha n$ for every $x$ \ploc{(8.16), p.~64}{terminal.tex:386--395}; split the curvature form over $S(x)$ \ploc{(p.~64)}{terminal.tex:397--421} & $\beta<1/2$: entropy contraction $\ge\rho_\beta/n$ for the polarized walk on Hamming balls of radius $\varepsilon(\beta)n$ (Thm~7.1, p.~110); additive error $e^{-Cn}$ (Thm~5.1, p.~21)\\
\addlinespace
\bottomrule
\end{xltabular}
\end{center}

\paragraph{How to check.}
\begin{enumerate}
\item The count \ploc{(p.~63)}{terminal.tex:370--384} against the proof of
\sloc{DaviesEtAl2026-PotentialHessianAscentIIISampling}{Lemma~7.43, p.~135}.
\item The $U_{SS}$ part of \ploc{(8.11)}{terminal.tex:235--239} and its proof
\ploc{(pp.~62--63)}{terminal.tex:297--314} against
\sloc{DaviesEtAl2026-PotentialHessianAscentIIISampling}{Lemma~7.23 and proof,
p.~115}.
\item The lead-in \ploc{(p.~61)}{terminal.tex:217--223} and the split
\ploc{(p.~64)}{terminal.tex:397--421} against PHA~III's overview (p.~19),
Theorem~7.1 and Remark~7.2 (p.~110); \ploc{(8.14)}{terminal.tex:342--348}
against \sloc{DaviesEtAl2026-WeakPoincareApproxSL}{Thm~5.1, p.~21}.
\end{enumerate}

\space%
\finding[Covariance at localization fields via the TAP Hessian]{A Poincar\'e inequality from covariance bounds at the localization fields, obtained from a uniformly positive TAP Hessian}{tap-hessian-covariance-route}

\paragraph{Summary.}
El Alaoui and Gaitonde
\sloc{ElAlaouiGaitonde2022-CovarianceMatrixSK}{p.~2 and \S2, pp.~3--4}
proposed a route to Poincar\'e and log-Sobolev inequalities at every
$\beta<1$: bound the covariance of the tilted measure $\mu_y$ in operator norm
at the fields of Eldan's stochastic localization, through its approximation by
the inverse negative TAP Hessian and a negative bound on that Hessian at the
mean, then transfer the bound along the localization following Chen--Eldan.
The preprint cites them only for their zero-field Theorem~1.1
\ploc{(\S1.1, p.~4)}{introduction.tex:89--92}, credits Chen--Eldan and
Celentano's local convexity but no work for the route joining them, and
presents its posterior estimates as ``additional'' to covariance results, which
``control linear fluctuations'' \ploc{(p.~4)}{introduction.tex:99--123}.
It carries out the route:
\ploc{Proposition~5.1 (p.~27)}{stability.tex:35--60} bounds $\mathsf H_J$
below, and $\mathsf H_J(y,V_y)$ is their negative TAP Hessian in normalized
coordinates; (1.5), proved as
\ploc{Proposition~7.2 (p.~56)}{consequences.tex:89--98}, is equivalent up to
constants to $\norm{\Cov_{\mu_h}(x)}\le C'$ at every good field, which the
preprint never states; and its \ploc{\S2 (pp.~7--10)}{observation.tex:1--210}
is the transfer, adding the square-reweighted mean (1.6) and a terminal gap.
The relationship is \rel{same route}: they state each step heuristically, and
the preprint proves each by its own means, the covariance bound by residual
identities tested against an implicit linear-response vector rather than
through their Frobenius bound (2.7); Brennecke, Xu and Yau
\sloc{BrenneckeXuYau2023-OperatorNormCorrelationSK}{Thm~1.4, p.~6}, also not
cited, had carried out that step at homogeneous fields $h\one$.

\paragraph{Correspondence.}
Their field $y$ is the preprint's $h$ ($h_0$ in its \S\S5--7); in the preprint
column, $y$ is the variable of (5.1)--(5.2) and $m=\tanh y$.
\begin{center}\small
\begin{xltabular}{\linewidth}{@{}>{\raggedright\arraybackslash}p{0.12\linewidth}>{\raggedright\arraybackslash}X>{\raggedright\arraybackslash}X@{}}
\toprule
Step & Preprint & El Alaoui--Gaitonde \cite{ElAlaouiGaitonde2022-CovarianceMatrixSK}\\
\midrule
\endfirsthead
\toprule
Step & Preprint & El Alaoui--Gaitonde \cite{ElAlaouiGaitonde2022-CovarianceMatrixSK}\\
\midrule
\endhead
Hessian & $\mathsf H_J(y,V_y)=V_y^{1/2}\Phi''(m)V_y^{1/2}$
\ploc{(proof of Lemma~5.7, p.~40)}{stability.tex:1143--1156} &
$-\nabla^2F_{\mathrm{TAP}}(m)$ at a maximizer $m$, (2.2)--(2.5), p.~3\\
\addlinespace
Margin & $\mathsf H_J(y,A)\succ cI$ for $\norm{F_{J,Y(t)}(y)}\le2\rho_0\sqrt n$,
$A$ near $V_y$, $t\le T$, off probability $e^{-an}$
\ploc{(Prop.~5.1, p.~27)}{stability.tex:35--60} &
$\lambda_{\max}(\nabla^2F_{\mathrm{TAP}}(m))\le-\varepsilon$ at the mean $m$ of
$\mu_y$, with probability $1-O(1/n)$ (p.~4)\\
\addlinespace
Covariance & (1.5), i.e.\ (7.2), at every field satisfying (5.3)
\ploc{(Prop.~7.2, p.~56)}{consequences.tex:89--98}; equivalently
$\norm{\Cov_{\mu_h}(x)}\le C'$ &
Bounded $\E\norm{\Cov(\mu_y)}_{\mathrm{op}}$, via the Frobenius form (2.7) of
$\Cov(\mu_y)\approx(-\nabla^2F_{\mathrm{TAP}}(m))^{-1}$ (p.~4)\\
\addlinespace
Transfer & $u'(t)=-\E[\one_{\{t<\tau\}}\norm{\Cov_t(f,x)}^2]$, bounded
through (1.5) \ploc{(\S2.2, p.~8)}{observation.tex:70--86} &
Bounds on $\norm{\Cov(\mu_y)}_{\mathrm{op}}$ at localization fields give
Poincar\'e and log-Sobolev for $\mu$ (p.~4)\\
\bottomrule
\end{xltabular}
\end{center}

\paragraph{How to check.}
\begin{enumerate}
\item \sloc{ElAlaouiGaitonde2022-CovarianceMatrixSK}{(2.3), (2.5), p.~3}
against \ploc{(5.2), p.~27}{stability.tex:22--33} and $\Phi$ in the proof of
\ploc{Lemma~5.7, p.~40}{stability.tex:1143--1156}: with $J=\beta A$,
$j=\beta^2$ and $D(m)=V_y^{-1}$,
$\Phi(m)=-F_{\mathrm{TAP}}(m)-h_0^{\mathsf T}m+n(\log2+j/4)$ and
$\mathsf H_J(y,V_y)=V_y^{1/2}(-\nabla^2F_{\mathrm{TAP}}(m))V_y^{1/2}$; as
$V_y\preceq I$, $\mathsf H_J\succeq cI$ gives
$\lambda_{\max}(\nabla^2F_{\mathrm{TAP}}(m))\le-c$.
\item $G=e^{\mathsf T}x$, $\norm e=1$, in
\ploc{(1.5), p.~6}{introduction.tex:212--217}: $\Cov_{\mu_h}(G,x)=\Cov_{\mu_h}(x)e$
and $\D_h(G)=\sum_ie_i^2\,\E\Var(x_i\mid x_{-i})\le1$, so
$s=\norm{\Cov_{\mu_h}(x)}$ has $s^2\le C(s+1)$, $s\le C+1$; Cauchy--Schwarz
gives the converse.
\item The two paragraphs after (2.7) in
\sloc{ElAlaouiGaitonde2022-CovarianceMatrixSK}{p.~4} against inputs~1--2 of
\ploc{\S1.4 (p.~6)}{introduction.tex:197--231} and \ploc{\S2}{observation.tex:1--210}.
\end{enumerate}

\space%
\finding[A weak Poincar\'e inequality for the SK heat-bath chain at \texorpdfstring{$\beta<1/2$}{beta<1/2}]{A weak Poincar\'e inequality for the zero-field SK heat-bath chain at \texorpdfstring{$\beta<1/2$}{beta<1/2}}{weak-poincare-beta-half}

\paragraph{Summary.}
Davies, Lee, Sandhu and Shi's weak Poincar\'e paper
\cite{DaviesEtAl2026-WeakPoincareApproxSL}, posted on July 9, 2026, proves
that for every $\beta<1/2$, with high probability over the disorder, the
zero-field SK Gibbs measure satisfies a weak Poincar\'e inequality for the
single-site heat-bath (Glauber) chain, with additive error
$\epsilon\operatorname{osc}(f)^2$ for every $\epsilon\in(0,1)$
\sloc{DaviesEtAl2026-WeakPoincareApproxSL}{Cor.~1.5, p.~2}. The preprint
does not cite this paper anywhere: its account of earlier results in \ploc{\S1.1 (pp.~3--4)}{introduction.tex:52--87} omits
it, although two later works cited there describe it, Wang's
\sloc{Wang2026-OptimalMixingGlauberSKBetaHalf}{p.~5} and Boban, Li and Oveis
Gharan's \sloc{BobanLiOveisGharan2026-RankOnePerturbedTrickledownSK}{p.~7}.
The earlier result is the inequality of the preprint's main result,
\ploc{Theorem~1.1 (p.~3)}{introduction.tex:33--42}, with an added error term
and only for $\beta<1/2$ instead of every $\beta<1$: \rel{weaker result},
since Theorem~1.1 restricted to $\beta<1/2$ implies it and not conversely.
Both proofs observe a spin sample in Gaussian noise up to a large fixed time,
prove an inequality for the posterior laws there and transfer it back, an
approach the preprint credits to Eldan and to Chen and Eldan in
\ploc{\S1.1 (p.~4)}{introduction.tex:105--110}; in the earlier paper,
exceptional localized laws and a weak inequality for the observation law
enter as additive $\operatorname{osc}(f)^2$ terms, while the preprint
weights paths by $f^2$ to bound the variance on exceptional ones
\ploc{(p.~4; \S2.4, p.~10)}{introduction.tex:111--113;
observation.tex:188--194}. The same paper is a technique source in
findings~\ref{f:stopped-localization}
and~\ref{f:terminal-gap-sparse-submatrices}.

\paragraph{Correspondence.}
\begin{center}\small
\begin{xltabular}{\linewidth}{@{}>{\raggedright\arraybackslash}p{0.13\linewidth}>{\raggedright\arraybackslash}X>{\raggedright\arraybackslash}X@{}}
\toprule
 & Preprint & Weak Poincar\'e paper \cite{DaviesEtAl2026-WeakPoincareApproxSL}\\
\midrule
\endfirsthead
\toprule
 & Preprint & Weak Poincar\'e paper \cite{DaviesEtAl2026-WeakPoincareApproxSL}\\
\midrule
\endhead
Measure & $\mu_0$; $J$ has zero diagonal and entries $J_{ij}\sim N(0,j/n)$, $j=\beta^2$ \ploc{(\S1, p.~3)}{introduction.tex:13--19} & $\mu_{\beta A}(\sigma)\propto e^{\frac\beta2\ip{\sigma}{A\sigma}}$, $A\sim\mathrm{GOE}(n)$ (p.~1): diagonal variance $2/n$, off-diagonal $1/n$ (PHA~III \sloc{DaviesEtAl2026-PotentialHessianAscentIIISampling}{p.~1})\\
\addlinespace
Dirichlet form & $\D_0(f)=\sum_i\E_{\mu_0}(f-P_i^0f)^2$ \ploc{(1.1), p.~3}{introduction.tex:22--24} & Glauber: $\mathcal E_\mu(f,f)=\frac1{2n}\sum_{x,i}\frac{\mu(x)\mu(x^{\oplus i})}{\mu(x)+\mu(x^{\oplus i})}\bigl(f(x^{\oplus i})-f(x)\bigr)^2$ (p.~6)\\
\addlinespace
Inequality & $\Var_{\mu_0}(f)\le C_\beta\D_0(f)$ for every $f$ \ploc{(Theorem~1.1, p.~3)}{introduction.tex:33--42} & $\Var_\mu(f)\le ne^{C/\epsilon}\mathcal E_\mu(f,f)+\epsilon\operatorname{osc}(f)^2$ for every $\epsilon\in(0,1)$: a $(\frac1ne^{-C/\epsilon},\epsilon)$-weak Poincar\'e inequality (Cor.~1.5, p.~2; Def.~2.3, p.~6, with $\mathrm{Err}=\operatorname{osc}^2$ as on p.~7)\\
\addlinespace
Range and disorder & every fixed $\beta<1$; probability tending to one; $C_\beta$ depends only on $\beta$ & $\beta<1/2$; probability at least $1-\delta$ for $n\ge N(\beta,\delta)$; $C=C(\beta,\delta)$\\
\bottomrule
\end{xltabular}
\end{center}

\paragraph{How to check.}
\begin{enumerate}
\item The Glauber Dirichlet form
\sloc{DaviesEtAl2026-WeakPoincareApproxSL}{p.~6} against
\ploc{(1.1), p.~3}{introduction.tex:22--24}: each pair $\{x,x^{\oplus i}\}$
differing in spin $i$ contributes its mass times the conditional variance of
$f$ on it, so $\mathcal E_{\mu_0}(f,f)=n^{-1}\D_0(f)$. The diagonal of $A$
adds a constant on $\{\pm1\}^n$ \ploc{(p.~3)}{introduction.tex:47--48}, so
$\mu_{\beta A}=\mu_0$ with $J$ the off-diagonal part of $\beta A$. Thus
Corollary~1.5 reads
$\Var_{\mu_0}(f)\le e^{C/\epsilon}\D_0(f)+\epsilon\operatorname{osc}(f)^2$
for $\beta<1/2$ and $0<\epsilon<1$.
\item \ploc{Theorem~1.1 (p.~3)}{introduction.tex:33--42} with $\beta<1/2$
against \sloc{DaviesEtAl2026-WeakPoincareApproxSL}{Cor.~1.5, p.~2}: the
first gives the second with any $C\ge\log C_\beta$, since
$e^{C/\epsilon}\ge e^{C}$ for $\epsilon<1$; the additive term blocks the
converse.
\item The preprint's \ploc{\S1.1 (pp.~3--4)}{introduction.tex:52--87} and
its bibliography, which do not mention the paper, against
\sloc{Wang2026-OptimalMixingGlauberSKBetaHalf}{p.~5} and
\sloc{BobanLiOveisGharan2026-RankOnePerturbedTrickledownSK}{p.~7}, which
describe it.
\end{enumerate}

\space%
\finding[Gaussian-comparison margin for the TAP Hessian]{A positive margin for the Onsager-shifted TAP Hessian by comparison with a norm-weighted Gaussian process}{sudakov-fernique-tap-margin}

\paragraph{Summary.}
The proof of \ploc{Lemma~3.2 (pp.~12--13)}{matrices.tex:145--236} names the
Sudakov--Fernique inequality and cites Vitale for adding the same
deterministic function to both processes \ploc{(p.~12)}{matrices.tex:169--173}.
The comparison it then makes, of the GOE quadratic form $v^{\mathsf T}Wv$
with the norm-weighted linear process $(2/\sqrt n)\norm v\ip gv$ under a
common deterministic shift, is Lemma~4.1 of Celentano, Fan and Mei
\sloc{CelentanoFanMei2021-LocalConvexityTAPZ2Sync}{Lemma~4.1, p.~15}, used
for the TAP Hessian in their Lemma~4.9 (p.~18; also Lemma~4.12); the
preprint does not cite that paper, and cites Celentano, who restates the
comparison \sloc{Celentano2022-SudakovFerniquePostAMPTAPConvexity}{Prop.~1, p.~6},
only for local convexity near AMP iterates
\ploc{(\S1.1, p.~4)}{introduction.tex:120--122}.
On \ploc{p.~13}{matrices.tex:179--221} the preprint takes $v\propto A^{1/2}p$
and the Onsager term as the shift, for each fixed diagonal $0\le A\le a_+I$
($a_+=1+\eta$), and bounds the smallest eigenvalue of $S_A(z;W)$ below through
a scalar supremum \ploc{(3.9)}{matrices.tex:197--217}; the resulting event
feeds \ploc{Lemmas~3.3--3.5}{matrices.tex:349--351, 716, 753--757} and
\ploc{5.5}{stability.tex:481--483, 638--639}.
The shared approach is Celentano, Fan and Mei's: to bound a TAP Hessian below
beyond what $\norm W$ gives, compare its random part with the norm-weighted
linear process while the deterministic terms, the Onsager term among them,
ride along as a common shift, reducing the smallest eigenvalue to a low-dimensional
variational problem; here this gives a margin on the whole range $ja_+^2<1$,
where $\norm W\le2\sqrt j+\varepsilon$ alone gives one only when
$2\sqrt j\,a_+<1$.
\rel{same technique}: the lemma is used intact, without conditioning and at a
fixed diagonal, where the variational problem becomes scalar; Celentano and,
not cited, Brennecke, Xu and Yau
\sloc{BrenneckeXuYau2023-OperatorNormCorrelationSK}{Prop.~4.5 and Step~1 of its proof, pp.~24--28}
build on the same use.

\paragraph{Correspondence.}
\begin{center}\small
\begin{xltabular}{\linewidth}{@{}>{\raggedright\arraybackslash}p{0.12\linewidth}>{\raggedright\arraybackslash}X>{\raggedright\arraybackslash}X@{}}
\toprule
 & Preprint & Celentano--Fan--Mei \cite{CelentanoFanMei2021-LocalConvexityTAPZ2Sync}\\
\midrule
\endfirsthead
\toprule
 & Preprint & Celentano--Fan--Mei \cite{CelentanoFanMei2021-LocalConvexityTAPZ2Sync}\\
\midrule
\endhead
Comparison & $\E\sup_{\norm p=1}\{zX_p-z^2j\chi_A\norm w^2\}\le\E\sup_{\norm p=1}\{zY_p-z^2j\chi_A\norm w^2\}$, $X_p=w^{\mathsf T}Ww$, $Y_p=2\sqrt{j/n}\,\norm w\,g^{\mathsf T}w$, $w=A^{1/2}p$ \ploc{(p.~13)}{matrices.tex:179--196} & $\E\sup_{x\in\mathcal X}[v(x)^{\mathsf T}Wv(x)+f(x)]\le\E\sup_{x\in\mathcal X}[\frac{2}{\sqrt n}\norm{v(x)}\ip{g}{v(x)}+f(x)]$, $g\sim N(0,I_n)$ (Lemma~4.1, p.~15)\\
\addlinespace
Increments & $\E(Y_p-Y_{p'})^2-\E(X_p-X_{p'})^2=\frac{2j}{n}\bigl[(\norm w^2-\norm{w'}^2)^2+2(\norm w\norm{w'}-w^{\mathsf T}w')^2\bigr]\ge0$ \ploc{(p.~13)}{matrices.tex:184--191} & $\E\ip v{Wv}\ip{v'}{Wv'}=\frac2n\ip v{v'}^2$; two inequalities raise $\frac2n(\norm v^4+\norm{v'}^4-2\ip v{v'}^2)$ to $\frac2n(2\norm v^4+2\norm{v'}^4-4\norm v\norm{v'}\ip v{v'})$; Sudakov--Fernique with common mean $f$ (App.~A.4, p.~33)\\
\bottomrule
\end{xltabular}
\end{center}

\paragraph{How to check.}
\begin{enumerate}
\item The comparison display on \ploc{p.~13}{matrices.tex:192--196} against
\sloc{CelentanoFanMei2021-LocalConvexityTAPZ2Sync}{Lemma~4.1, p.~15}, with
$\mathcal X$ the unit sphere, $v(p)=(z\sqrt j)^{1/2}A^{1/2}p$,
$f(p)=-z^2j\chi_Ap^{\mathsf T}Ap$ and the preprint's $W$ equal to $\sqrt j$
times $\mathrm{GOE}(n)$: their two processes become $zX_p$ and $zY_p$.
\item The increment identity on \ploc{p.~13}{matrices.tex:184--191} against
\sloc{CelentanoFanMei2021-LocalConvexityTAPZ2Sync}{App.~A.4, p.~33}: the
preprint's bracket is the sum of the two terms the two inequalities drop.
\item $S_A(1;W)=A^{1/2}(A^{-1}+j\chi_AI-W)A^{1/2}$
\ploc{(\S3, p.~11)}{matrices.tex:34--37} against the TAP Hessian $H(m)$ of
\sloc{CelentanoFanMei2021-LocalConvexityTAPZ2Sync}{(4.16), p.~16}: with
$j=\lambda^2$, $A=\operatorname{diag}(1-m^2)$ and $\chi_A=1-\norm m^2/n$,
$A^{-1}+j\chi_AI-W$ is $H(m)$ without the spike
$-\frac{\lambda^2}{n}\one\one^{\mathsf T}$ and the term
$-\frac{2\lambda^2}{n}mm^{\mathsf T}$, and the Onsager term is deterministic
in both.
\end{enumerate}

\space%
\finding[Consistent Gaussian field law on the Nishimori line]{The consistent Gaussian field law: mean equals variance, then a one-variable fixed point}{nishimori-consistent-gaussian}

\paragraph{Summary.}
El Alaoui, Montanari and Sellke (EAMS below) reduce the state evolution of
AMP in the planted SK model with Gaussian side observation $tx+\sqrt t\,W$ to
the fixed point $\gamma_*=\beta^2(1-\operatorname{mmse}(\gamma_*+t))$, and
prove its existence, its uniqueness for $t>0$ (by concavity), its size of
order $t$ and its Lipschitz dependence on $t$ (the last two by comparing
fixed-point equations with slope at most $\beta^2$)
\sloc{ElAlaouiMontanariSellke2022-SamplingSKAlgorithmicSL}{Prop.~4.4, (4.20)--(4.24), pp.~13--15};
for $\beta^2<1/3$, Davies, Lee, Sandhu and Shi's Potential Hessian Ascent III
(PHA III) proves that every solution of the two-parameter system behind it
has equal mean and variance
\sloc{DaviesEtAl2026-PotentialHessianAscentIIISampling}{Lemma~5.32, pp.~98--99}.
The preprint cites EAMS only in \ploc{\S1.1, p.~4}{introduction.tex:118--120},
for their sampling method, and does not cite PHA III; its \S4 cites nothing
and does not name the Nishimori identity
$\E_{N(s,s)}\tanh y=\E_{N(s,s)}\tanh^2y$ \ploc{(4.10)}{entropy.tex:250--259}
on which both steps rest.
In \ploc{Lemma~4.2 (pp.~25--26)}{entropy.tex:220--337} it solves the
consistency equations $d=t+j\E_P\tanh y$, $s=t+j\E_P\tanh^2y$ of
\ploc{(4.7)}{entropy.tex:206--218} (with $\sigma=0$) for Gaussian laws
$P=N(d,s)$: it forces $d=s$ by the same Lipschitz comparison as PHA III, with
constant $1$ on $[0,\infty)$ in place of $3$, so for every $j<1$; the
remaining equation \ploc{(4.8)}{entropy.tex:224--228} is EAMS's (4.21) under
$s_*=t+\gamma_*$, $j=\beta^2$, with the same existence, uniqueness, size and
Lipschitz statements and the same comparison for the Lipschitz bound.
This is \rel{same technique}: the same scalar system, fixed point and bounds,
with the step $d=s$ as in PHA III's Lemma~5.32; the preprint proves
uniqueness by a contraction rather than by concavity, and includes $t=0$.
EAMS credit the one-variable reduction and the properties of mmse to
Deshpande, Abbe and Montanari
\cite{DeshpandeAbbeMontanari2015-AsymptoticMutualInfoTwoGroupsSBM}, who
use erasure side information in place of the Gaussian observation.

\paragraph{Correspondence.}
With $j=\beta^2$, $s_*=t+\gamma_*$ and, by (4.10),
$\E_{N(s,s)}\tanh y=1-\operatorname{mmse}(s)$ (EAMS), and with
$\alpha=t+\beta^2q^*=s$, $\delta=\beta^2(m^*-q^*)=d-s$ (PHA III):
\begin{center}\small
\begin{xltabular}{\linewidth}{@{}>{\raggedright\arraybackslash}p{0.11\linewidth}>{\raggedright\arraybackslash}X>{\raggedright\arraybackslash}X@{}}
\toprule
 & Preprint & EAMS \cite{ElAlaouiMontanariSellke2022-SamplingSKAlgorithmicSL}; PHA III \cite{DaviesEtAl2026-PotentialHessianAscentIIISampling} where named\\
\midrule
\endfirsthead
\toprule
 & Preprint & EAMS \cite{ElAlaouiMontanariSellke2022-SamplingSKAlgorithmicSL}; PHA III \cite{DaviesEtAl2026-PotentialHessianAscentIIISampling} where named\\
\midrule
\endhead
$d=s$ & $H_s(d)=\E_{N(d,s)}[\tanh y-\tanh^2y]$, 1-Lipschitz on $[0,\infty)$, $H_s(s)=0$ by (4.10); $|H_s(d)|\le|d-s|=j|H_s(d)|$ \ploc{(pp.~25--26)}{entropy.tex:269--298} & PHA III: $m^*-q^*=\E h(\alpha+\delta+\sqrt\alpha G)$, $h=\tanh-\tanh^2$ (5.81), $\E h(\alpha+\sqrt\alpha G)=0$ (5.82) (p.~98); $|h'|\le3$, $|m^*-q^*|\le3|\delta|=3\beta^2|m^*-q^*|$ (p.~99)\\
\addlinespace
Fixed point & \ploc{(4.8)}{entropy.tex:224--228}: $s_*=t+j\E_{N(s_*,s_*)}\tanh y$; $q_*=(s_*-t)/j$ \ploc{(4.9), p.~25}{entropy.tex:236--241} & (4.21): $\gamma_*=\beta^2(1-\operatorname{mmse}(\gamma_*+t))$ (p.~14); $q_*=\gamma_*/\beta^2$, the limit of (2.5) (p.~5)\\
\addlinespace
Size & \ploc{(4.9)}{entropy.tex:236--241}: $t\le s_*\le t/(1-j)$ & (4.23); its proof gives $\gamma_*\le\beta^2t/(1-\beta^2)$ (p.~15)\\
\addlinespace
Lipschitz & $|s_*(t)-s_*(t')|\le|t-t'|+j|s_*(t)-s_*(t')|$ \ploc{(p.~26)}{entropy.tex:323--329} & Proof of (4.24): $|\gamma_*(t_1)-\gamma_*(t_2)|\le\beta^2|\gamma_*(t_1)-\gamma_*(t_2)|+\beta^2|t_1-t_2|$ (p.~15)\\
\bottomrule
\end{xltabular}
\end{center}

\paragraph{How to check.}
\begin{enumerate}
\item The $d=s$ argument \ploc{(pp.~25--26)}{entropy.tex:269--298} against
\sloc{DaviesEtAl2026-PotentialHessianAscentIIISampling}{Lemma~5.32 and proof, pp.~98--99}:
$H_s(d)$ with (5.81), (4.10) with (5.82), $j$ with $3\beta^2$.
\item \ploc{(4.8)--(4.9), p.~25}{entropy.tex:224--241} and the Lipschitz
comparison on \ploc{p.~26}{entropy.tex:323--329} against
\sloc{ElAlaouiMontanariSellke2022-SamplingSKAlgorithmicSL}{(4.21), (4.23)--(4.24), p.~14}
and their proofs (p.~15), with $s_*=t+\gamma_*$ and $j=\beta^2$.
\end{enumerate}

\space%
\finding[Reweighted path entropy as entropy decay]{The entropy of the reweighted observation law as the entropy lost along the localization}{tilted-path-entropy}

\paragraph{Summary.}
The preprint credits Lehec, and F\"ollmer through Lehec
\cite{Lehec2010-RepresentationFormulaEntropy}, for the drift-energy identity
\ploc{(2.5)}{observation.tex:107--112}, which writes the relative entropy of the stopped observation path laws
under $f^2\mu_0$ and under $\mu_0$ as half the expected drift energy
(\ploc{\S1.1, p.~4}{introduction.tex:114--116}, and at the point of use,
\ploc{\S2.3, p.~8}{observation.tex:104--105}). It does not cite Chen and Eldan
for what this quantity is along the localization: the entropy of $f^2$ lost,
at the rate given by their entropy-decay formula
\sloc{ChenEldan2022-LocalizationSchemes}{\S3.2.3, eq.~(27), p.~29}; it cites
them only for their variance-retention and Dirichlet-form framework
(\ploc{\S1.1, p.~4}{introduction.tex:108--110};
\ploc{\S2, p.~7}{observation.tex:8--9}). In
\ploc{\S2.3 (pp.~8--9)}{observation.tex:90--160}, (2.5) is equivalent to the
expectation of (27) with $C_t=I_n$ and $f^2$ in place of $f$, integrated up to
a stopping time $\tau\le T$, and \ploc{(2.7)}{observation.tex:130--136} bounds
this entropy so as to make the bad paths rare under the reweighted law in
\ploc{(2.8)--(2.9)}{observation.tex:139--160}, a step toward
\ploc{Theorem~1.1}{introduction.tex:33--42}. \rel{fragment}: the quantity that
(2.5) expresses and (2.7) bounds is the entropy loss in (27), but the preprint
proves (2.5) by It\^o's formula for $\log N_t$, as in the proof of Lehec's
Thm~2 \sloc{Lehec2010-RepresentationFormulaEntropy}{pp.~4--5}, bounds the mean
gap $a_t$ by comparing both means with the TAP root through
\ploc{(1.6)}{introduction.tex:218--224} instead of by Chen and Eldan's
entropic stability \sloc{ChenEldan2022-LocalizationSchemes}{Def.~29, p.~25},
and has no counterpart of their Gr\"onwall step and entropy conservation
\sloc{ChenEldan2022-LocalizationSchemes}{Prop.~39, p.~29}.

\paragraph{Correspondence.}
Chen and Eldan's $\nu_t$ (the localized measure) and $\mu_t$ (its reweighting)
swap the preprint's letters.
\begin{center}\small
\begin{xltabular}{\linewidth}{@{}>{\raggedright\arraybackslash}p{0.15\linewidth}>{\raggedright\arraybackslash}X>{\raggedright\arraybackslash}X@{}}
\toprule
 & Preprint & Chen--Eldan \cite{ChenEldan2022-LocalizationSchemes}\\
\midrule
\endfirsthead
\toprule
 & Preprint & Chen--Eldan \cite{ChenEldan2022-LocalizationSchemes}\\
\midrule
\endhead
Localization & posterior $\mu_{Y(t)}$ of $X\sim\mu_0$, $Y(t)=tX+B(t)$, \ploc{(1.4), p.~6}{introduction.tex:199--201}; innovation $I(t)$, \ploc{\S2.1, p.~7}{observation.tex:12--21} & $\nu_t$ from (5) with $C_t=I_n$, the Gaussian channel, driven by $B_t$ (pp.~16--17)\\
\addlinespace
Reweighting & $\nu_t=f^2\mu_{Y(t)}/N_t$, $N_t=\E_tf^2$, $\E_{\mu_0}f^2=1$; $\P^f$, $\E^f$: observation law for $X\sim f^2\mu_0$, density $N_t$, \ploc{p.~8}{observation.tex:90--97} & $d\mu_t/d\nu_t=f/\nu_t(f)$, $f\ge0$; $M_t=\nu_t(f)$ (p.~28)\\
\addlinespace
Mass martingale & $dN_t/N_t=a_t^{\mathsf T}dI(t)$, $a_t=\E_{\nu_t}x-\E_tx$, \ploc{p.~8}{observation.tex:98--103} & $dM_t=M_t\ip{C_t(\mathrm b(\mu_t)-\mathrm b(\nu_t))}{dB_t}$, $\mathrm b$ the barycenter (p.~28)\\
\addlinespace
Entropy & $D(\P^f|_{\mathcal F_\tau}\Vert\P|_{\mathcal F_\tau})=\E^f\log N_\tau=\tfrac12\E^f\!\int_0^\tau\norm{a_t}^2dt$, It\^o for $\log N_t$, \ploc{(2.5), p.~8}{observation.tex:104--116} & $d\Ent_{\nu_t}[f]=-\tfrac12\nu_t(f)\,|C_t(\mathrm b(\mu_t)-\mathrm b(\nu_t))|^2dt+{}$martingale, It\^o for $M_t\log M_t$, (27), p.~29\\
\bottomrule
\end{xltabular}
\end{center}
Since $\E_t[f^2\log f^2]$ is a $\P$-martingale and $N_0=1$,
$\E^f\log N_\tau=\Ent_{\mu_0}(f^2)-\E\Ent_{\mu_{Y(\tau)}}(f^2)$ and
$\E^f\int_0^\tau\norm{a_t}^2dt=\E\int_0^\tau N_t\norm{a_t}^2dt$, so (2.5) is
equivalent to the expectation of (27), with $C_t=I_n$ and $f$ replaced by
$f^2$, integrated from $0$ to $\tau$.

\paragraph{How to check.}
\begin{enumerate}
\item The formula for $dN_t/N_t$ and \ploc{(2.5) (p.~8)}{observation.tex:98--116}
against $dM_t$ and (27) in
\sloc{ChenEldan2022-LocalizationSchemes}{\S3.2.3, pp.~28--29}, with $C_t=I_n$
(p.~17) and $f^2$ in place of $f$, through the martingale
$\E_t[f^2\log f^2]$.
\item The proof of (2.5) by It\^o's formula for $\log N_t$
\ploc{(p.~8)}{observation.tex:104--116} against the proof of
\sloc{Lehec2010-RepresentationFormulaEntropy}{Thm~2, pp.~4--5}.
\item The bound on $\norm{a_t}/\sqrt n$ from (1.6) and the bound (2.7)
\ploc{(p.~8)}{observation.tex:118--136} against Def.~29 and Lemma~31 (p.~25)
and \sloc{ChenEldan2022-LocalizationSchemes}{Prop.~39, p.~29}.
\end{enumerate}

\space%
\finding[Mixture-Gaussian Stein operator for tanh differences]{Comparing \texorpdfstring{$\tanh$}{tanh} at two fields through the Stein operator of a two-Gaussian mixture}{mixture-gaussian-stein-kernel}

\paragraph{Summary.}
Chatterjee \sloc{Chatterjee2007-SpinGlassesSteinsMethod}{\S1.8 and Lemma~1.3, pp.~7--8}
develops ``Stein's method for mixture Gaussians''. For the law
$M(a,b,\mu,\sigma^2)$, with density proportional to
$\cosh(ax+b)e^{-(x-\mu)^2/2\sigma^2}$, a two-Gaussian mixture under which
$\tanh(ax+b)$ has mean $\tanh(a\mu+b)$, the paper gives a characterizing operator and
bounded solutions of its Stein equation, and uses them, centred at the TAP
field, to explain and prove the TAP equations. The preprint does not cite this
paper, and its compiled text does not mention Stein's method. For $a>0$, its
scalar identity \ploc{(6.9), p.~45}{residuals.tex:255--260} is $-a$ times
Chatterjee's Stein equation with $u=\tanh$, for the parameters $a=1$, $b=0$, $\mu=r_0$
and $\sigma^2$ equal to the preprint's $a$, and its kernel $\Phi$
\ploc{(6.8)}{residuals.tex:249--254} is $-1/a$ times the bounded solution there.
The shared approach, and Chatterjee's insight (Chatterjee's \S1.4, p.~4, is titled
``Onsager correction and mixture Gaussians''), is to write $\tanh$ at a field
minus $\tanh$ at a centre as this operator, whose multiplicative part is the
field minus the centre minus a coefficient times $\tanh$ of the field, applied
to a bounded function, leaving terms that integration by parts controls; the
preprint does this in each induction step of its weak residual rule
\ploc{(Lemma~6.5, pp.~53--54)}{residuals.tex:886--1032} and, centred at the
root of the field equation, in
\ploc{(6.12)--(6.13), pp.~45--46,}{residuals.tex:316--342} and
\ploc{(7.3), p.~57}{consequences.tex:113--132}. \rel{fragment}: the preprint
uses only the operator and its $\tanh$ solution, as an exact identity whose
$a$ is a difference of Onsager coefficients of either sign, without the
mixture law, Lemma~1.3 for general $u$, or Gaussian integration by parts in
the disorder.

\paragraph{Correspondence.}
\begin{center}\small
\begin{xltabular}{\linewidth}{@{}>{\raggedright\arraybackslash}p{0.13\linewidth}>{\raggedright\arraybackslash}X>{\raggedright\arraybackslash}X@{}}
\toprule
 & Preprint & Chatterjee \cite{Chatterjee2007-SpinGlassesSteinsMethod}\\
\midrule
\endfirsthead
\toprule
 & Preprint & Chatterjee \cite{Chatterjee2007-SpinGlassesSteinsMethod}\\
\midrule
\endhead
Operator & $(z-r_0-a\tanh z)-a\partial_z$, $z,r_0,a\in\R$ \ploc{(6.9)}{residuals.tex:255--260} & $Tf=f'-\bigl(\frac{x-\mu}{\sigma^2}-a\tanh(ax+b)\bigr)f$ for $M(a,b,\mu,\sigma^2)$, $\sigma^2>0$ (12), p.~7\\
\addlinespace
Right side & $\tanh z-\tanh r_0$ \ploc{(6.9)}{residuals.tex:255--260} & $u-\int u\,\psi$ (15), equal to $\tanh x-\tanh\mu$ for $u=\tanh$ by (14), p.~7\\
\addlinespace
Bounded solution & $\Phi$ \ploc{(6.8)}{residuals.tex:249--254}; $0<\Phi\le e^{|a|/2}$; for $a$ in a bounded interval, fixed-order derivatives at most $C\Phi$, uniformly in $z,r_0$ \ploc{(p.~45)}{residuals.tex:265--273} & $f=\rho^{-1}\int_{-\infty}^x\rho\,(u-r(\mu))$, display before (54), p.~25; $|f|,|\partial_xf|,|\partial_\mu f|\le C(a,\sigma)\norm{u}_\infty$ (Lemma~1.3, p.~8)\\
\addlinespace
Centre & $r_0=h_i^{p+1}$ \ploc{(Lemma~6.5, p.~54)}{residuals.tex:966--972}; the root $r_i$ \ploc{(6.12)}{residuals.tex:316--320} & the TAP field $r_1$ of (10), p.~5; law $M(\beta,h,r_1,1-q)$ (pp.~20--22)\\
\addlinespace
\bottomrule
\end{xltabular}
\end{center}

\paragraph{How to check.}
\begin{enumerate}
\item \ploc{(6.9), p.~45}{residuals.tex:255--260} against $T$ and (14)--(15)
in \sloc{Chatterjee2007-SpinGlassesSteinsMethod}{\S1.8, p.~7}, with
Chatterjee's $a=1$, $b=0$, $\mu=r_0$ and $\sigma^2$ equal to the preprint's
$a>0$: $-\sigma^2T$ is the preprint's operator, and for $u=\tanh$ the right
side of (15) is $\tanh z-\tanh r_0$.
\item \ploc{(6.8) and the bounds below it, p.~45}{residuals.tex:249--274}
against the display before (54) in
\sloc{Chatterjee2007-SpinGlassesSteinsMethod}{\S3.6, p.~25} with $u=\tanh$,
and Lemma~1.3 (p.~8): for $a>0$, $-a\Phi$ is a bounded solution of (15), and
the other solutions differ by multiples of
$1/\rho=e^{(x-\mu)^2/2\sigma^2}/\cosh x$, which is unbounded, so
$\Phi=-f/a$.
\item \ploc{(6.12)--(6.13), pp.~45--46,}{residuals.tex:316--342} and
\ploc{(7.3), p.~57}{consequences.tex:113--132} against the centring at the TAP
field $r_1$ in
$\langle(\ell_1-r_1-\beta(1-q)\tanh(\beta\ell_1+h))f\rangle-(1-q)\langle f_1\rangle=-(1-q)\langle Tf\rangle$,
with $\ell_1$ the local field
\sloc{Chatterjee2007-SpinGlassesSteinsMethod}{p.~22}.
\end{enumerate}

\space%
\finding[Posterior mean and TAP root via a late iterate]{Comparing a posterior mean with the TAP root via a late iterate}{posterior-mean-tap-root}

\paragraph{Summary.}
El Alaoui, Montanari and Sellke, for $\beta<1/2$, place the mean $m(A,y)$ of
the SK Gibbs measure at a stochastic-localization field within
$\rho\sqrt{tn}$ of a stationary point $m_*$ of the TAP free energy through a
late AMP iterate: the mean is close to it, its small increments make it an
approximate TAP solution, and local strong convexity puts $m_*$ near it
\sloc{ElAlaouiMontanariSellke2022-SamplingSKAlgorithmicSL}{Lemma~4.10(2), p.~20; proof, p.~40}.
The preprint cites that paper only for its sampling method, in
\ploc{\S1.1 (p.~4)}{introduction.tex:118--120}, and cites nothing in \ploc{\S6.4 (p.~45)}{residuals.tex:288--314} or
\ploc{Proposition~7.3 (p.~58)}{consequences.tex:198--243}, where it compares
the mean of $\nu=f^2\mu_{h_0}/N$, $N=\E_{\mu_{h_0}}f^2$, with $t_*$ through a
late iterate of its recursion; both papers need the comparison for the drift
of a localization process, in
\sloc{ElAlaouiMontanariSellke2022-SamplingSKAlgorithmicSL}{Lemma~4.14, p.~23}
and in \ploc{\S2.3 (p.~8)}{observation.tex:118--123}. It
cites Bolthausen for its recursion; Bolthausen's iterates converge
\sloc{Bolthausen2012-IterativeTAPConstruction}{Thm~2.1, p.~4} but are not
compared with the Gibbs measure
\sloc{Bolthausen2012-IterativeTAPConstruction}{p.~2}.
\rel{fragment}: the preprint shares the outline and replaces both substantive
links: for a recursion started at a Gibbs sample and under $f^2$
reweighting, its mean--iterate link is an $O(1)$ bound in expectation
\ploc{(Lemma~6.5)}{residuals.tex:886--906} built by induction on the exact
conditional-spin identity \ploc{(6.31), p.~53}{residuals.tex:919--923}, with
small increments from pigeonhole \ploc{(Lemma~7.1)}{consequences.tex:18--83};
its root link is the global bound of
\ploc{Lemma~5.7 (p.~40)}{stability.tex:1105--1125}, not convexity near an
iterate.
The outline appears earlier in Celentano, Fan and Mei for
$\mathbb Z_2$-synchronization, from a Bayes-optimal late AMP iterate to a TAP
local minimizer by local strong convexity
\sloc{CelentanoFanMei2021-LocalConvexityTAPZ2Sync}{Thm~2.1(a), p.~8; p.~18}
(not cited); Celentano's convexity theorem, which the preprint cites in
\ploc{\S1.1}{introduction.tex:119--121}, extends El Alaoui,
Montanari and Sellke's link to all $\beta<1$
\sloc{Celentano2022-SudakovFerniquePostAMPTAPConvexity}{Thm~2, pp.~10--11; p.~4}.

\paragraph{Correspondence.}
\begin{center}\small
\begin{xltabular}{\linewidth}{@{}>{\raggedright\arraybackslash}p{0.13\linewidth}>{\raggedright\arraybackslash}X>{\raggedright\arraybackslash}X@{}}
\toprule
 & Preprint & El Alaoui--Montanari--Sellke\\
\midrule
\endfirsthead
\toprule
 & Preprint & El Alaoui--Montanari--Sellke\\
\midrule
\endhead
Mean $\approx$ iterate & $\lvert\E_\nu[\widehat C_k(x-m^k)^{\mathsf T}e]\rvert\le C_\delta(1+\D_{h_0}(f)/N)$, $\norm e\le1$, cutoffs $\widehat C_k$ of Lemma~7.1 (proof of Prop.~7.3, p.~58) & $\norm{m(A,y)-\hat m^k}/\norm{m(A,y)}\to0$, AMP iterate $\hat m^k=\tanh z^k$ of (4.15) (Prop.~4.8, (4.37), p.~17)\\
\addlinespace
Iterate $\approx$ approx.\ root & $F_{J,h_0}(h^k)=JR_k-jb_{k-1}(R_{k-1}+R_k)+j(b_k-b_{k-1})m^k$, $\lvert b_k-b_{k-1}\rvert\le2\norm{R_k}/\sqrt n$ (p.~45) & $\norm{z^{k+1}-z^k}/\norm{z^k}\to0$ (4.38) gives $\norm{\nabla\mathcal F_{\mathrm{TAP}}(\hat m^k;y,q)}/\sqrt{tn}\to0$ (Lemma~4.11, p.~21)\\
\addlinespace
Approx.\ root $\approx$ root & $\norm{y-r}\le C\norm{F_{J,h_0}(y)}$ if $\norm{F_{J,h_0}(y)}<r_0\sqrt n$, given $\mathsf H_J(y,V_y)\succeq cI$ on that whole region ((5.29), p.~40) & convexity for $\beta\le1/2-c$ (Lemma~4.12, p.~22) and Lemma~A.6 (p.~35): $\norm{\hat m^{K_{\mathsf{AMP}}}-m_*}\le\rho\sqrt{tn}/2$ (p.~40)\\
\addlinespace
Conclusion & $\norm{\E_\nu x-t_*}/\sqrt n\le C\delta+C_\delta n^{-1/2}(1+\D_{h_0}(f)/N)$ ((7.6), p.~58) & $\norm{m(A,y)-m_*(A,y)}_2\le\rho\sqrt{tn}$ (Lemma~4.10(2), p.~20)\\
\bottomrule
\end{xltabular}
\end{center}

\paragraph{How to check.}
\begin{enumerate}
\item The proof of \ploc{Proposition~7.3, p.~58}{consequences.tex:214--243}
against the proof of
\sloc{ElAlaouiMontanariSellke2022-SamplingSKAlgorithmicSL}{Lemma~4.10, Part~2, p.~40}
and the argument for \sloc{CelentanoFanMei2021-LocalConvexityTAPZ2Sync}{Thm~2.1(a), p.~18}:
mean to late iterate, iterate to root, triangle.
\item The $F_{J,h_0}(h^k)$ identity after \ploc{(6.11), p.~45}{residuals.tex:306--314}
against the proof of
\sloc{ElAlaouiMontanariSellke2022-SamplingSKAlgorithmicSL}{Lemma~4.11, p.~21},
with $\beta A\to J$, $\beta^2\to j$, $y\to h_0$, $z^k\to h^k$,
$q=\norm{\hat m^k}^2/n$.
\end{enumerate}

\space%

\begin{thebibliography}{99}
\small
\bibitem{ElAlaouiMontanariSellke2022-SamplingSKAlgorithmicSL} A. El Alaoui, A. Montanari, M. Sellke, "Sampling from the Sherrington-Kirkpatrick Gibbs measure via algorithmic stochastic localization", FOCS 2022 (full version). arXiv:2203.05093v2.
\bibitem{NL-AizenmanLebowitzRuelle1987} M. Aizenman, J. L. Lebowitz, D. Ruelle, Some rigorous results on the Sherrington--Kirkpatrick spin glass model, Comm. Math. Phys. 112 (1987) 3--20, doi:10.1007/BF01217677
\bibitem{HuangEtAl2024-WeakPoincareSimulatedAnnealing} B. Huang, S. Mohanty, A. Rajaraman, D. X. Wu, "Weak Poincaré inequalities, simulated annealing, and sampling from spherical spin glasses" (Nov 2024; v2). arXiv:2411.09075v2.
\bibitem{ChenEldan2022-LocalizationSchemes} Y. Chen, R. Eldan, "Localization schemes: A framework for proving mixing bounds for Markov chains", Duke Math. J. 174 (2025). arXiv:2203.04163v2.
\bibitem{JekelSandhuShi2024-PotentialHessianAscentSK} D. Jekel, J. S. Sandhu, J. Shi, "Potential Hessian Ascent: The Sherrington-Kirkpatrick Model" (Aug 2024; v3 May 20, 2026). arXiv:2408.02360v3.
\bibitem{DaviesEtAl2026-PotentialHessianAscentIIISampling} E. Davies, H. Lee, J. S. Sandhu, J. Shi, "Potential Hessian Ascent III: Sampling the Sherrington-Kirkpatrick model at beta < 1/2" (May 5, 2026; v2 May 19, 2026). arXiv:2605.03718v2.
\bibitem{FanMeiMontanari2018-TAPFreeEnergyVariationalInference} Z. Fan, S. Mei, A. Montanari, "TAP free energy, spin glasses, and variational inference", Ann. Probab. 49 (2021). arXiv:1808.07890v2.
\bibitem{Speicher1998-OperatorValuedFreeProbability} R. Speicher, "Combinatorial theory of the free product with amalgamation and operator-valued free probability theory", Mem. Amer. Math. Soc. 132 (1998), no. 627. doi:10.1090/memo/0627.
\bibitem{HuangSellke2023-StrongTopologicalTrivialization} B. Huang, M. Sellke, "Strong topological trivialization of multi-species spherical spin glasses" (Aug 2023; v3 May 2025). arXiv:2308.09677v3.
\bibitem{HuangMontanariPham2024-SphericalSpinGlassSamplingSL} B. Huang, A. Montanari, H. T. Pham, "Sampling from spherical spin glasses in total variation via algorithmic stochastic localization" (2024). arXiv:2404.15651v1.
\bibitem{CelentanoFanMei2021-LocalConvexityTAPZ2Sync} M. Celentano, Z. Fan, S. Mei, "Local convexity of the TAP free energy and AMP convergence for Z2-synchronization", Ann. Statist. 51 (2023). arXiv:2106.11428v3.
\bibitem{BayatiLelargeMontanari2012-UniversalityMessagePassing} M. Bayati, M. Lelarge, A. Montanari, "Universality in polytope phase transitions and message passing algorithms", Ann. Appl. Probab. 25 (2015) 753-822. arXiv:1207.7321v2; doi:10.1214/14-AAP1010.
\bibitem{Montanari2018-OptimizationSKHamiltonian} A. Montanari, "Optimization of the Sherrington-Kirkpatrick Hamiltonian" (Dec 2018; v2 Apr 2019). arXiv:1812.10897v2.
\bibitem{DaviesEtAl2026-WeakPoincareApproxSL} E. Davies, H. Lee, J. S. Sandhu, J. Shi, "Weak Poincaré inequalities via approximate stochastic localization: application to sampling the Sherrington-Kirkpatrick model" (July 9, 2026). arXiv:2607.08160v1.
\bibitem{BandeiraElAlaouiRodder2026-GlauberHighOverlapGibbs} A. S. Bandeira, A. El Alaoui, A. Rödder, "Mixing of Glauber dynamics on high overlap Gibbs measures" (July 2026). arXiv:2607.06813v1.
\bibitem{ElAlaouiGaitonde2022-CovarianceMatrixSK} A. El Alaoui, J. Gaitonde, "Bounds on the covariance matrix of the Sherrington-Kirkpatrick model" (Dec 2022; v3 Nov 2024). arXiv:2212.02445v3.
\bibitem{BrenneckeXuYau2023-OperatorNormCorrelationSK} C. Brennecke, C. Xu, H.-T. Yau, "Operator norm bounds on the correlation matrix of the SK model at high temperature" (2023; v2 Jan 2025). arXiv:2307.12535v2.
\bibitem{Celentano2022-SudakovFerniquePostAMPTAPConvexity} M. Celentano, "Sudakov-Fernique post-AMP, and a new proof of the local convexity of the TAP free energy", Ann. Probab. 52 (2024). arXiv:2208.09550v1.
\bibitem{Lehec2010-RepresentationFormulaEntropy} J. Lehec, "Representation formula for the entropy and functional inequalities", Ann. Inst. Henri Poincaré Probab. Stat. 49(3) (2013) 885-899. arXiv:1006.3028v2; doi:10.1214/11-AIHP464.
\bibitem{Chatterjee2007-SpinGlassesSteinsMethod} S. Chatterjee, "Spin glasses and Stein's method" (2007; v3 2009). arXiv:0706.3500v3.
\bibitem{OAI-Preprint} OpenAI, \emph{A spectral gap throughout the high-temperature Sherrington--Kirkpatrick phase}, openai/math release preprint, September 24, 2026, \url{https://github.com/openai/math}.
\bibitem{Eldan2012-ThinShellSpectralGapStochLoc} R. Eldan, "Thin shell implies spectral gap up to polylog via a stochastic localization scheme", GAFA 23 (2013). arXiv:1203.0893v5.
\bibitem{ElAlaouiEtAl2023-FastRelaxationRFIM} A. El Alaoui, R. Eldan, R. Gheissari, A. Piana, "Fast relaxation of the random field Ising dynamics" (Nov 10, 2023; v2 Nov 12, 2024). arXiv:2311.06171v2.
\bibitem{BeliusConcettiGenovese2024-DeterminantBrayMooreTAP} D. Belius, F. Concetti, G. Genovese, "On the determinant in Bray-Moore's TAP complexity formula" (2024). arXiv:2401.08529v1.
\bibitem{Boursier2026-LegendreStructureTAPComplexity} J. Boursier, "The Legendre structure of the TAP complexity for the Ising spin glass" (Apr 2026). arXiv:2604.20660v2.
\bibitem{Bolthausen2012-IterativeTAPConstruction} E. Bolthausen, "An iterative construction of solutions of the TAP equations for the Sherrington-Kirkpatrick model", Commun. Math. Phys. 325 (2014) 333-366. arXiv:1201.2891v1.
\bibitem{Wang2026-OptimalMixingGlauberSKBetaHalf} S. Wang, "Optimal mixing of Glauber dynamics for the Sherrington-Kirkpatrick model at beta<1/2" (v1 Aug 23, 2026; v2 Aug 28, 2026). arXiv:2608.22159v2.
\bibitem{BobanLiOveisGharan2026-RankOnePerturbedTrickledownSK} M. Boban, A. Li, S. Oveis Gharan, "Rank-1-perturbed trickledown theorems: Mixing time of Glauber dynamics for the Sherrington-Kirkpatrick model up to beta<=1/2+eps" (Sept 11, 2026). arXiv:2609.13138v1.
\bibitem{DeshpandeAbbeMontanari2015-AsymptoticMutualInfoTwoGroupsSBM} Y. Deshpande, E. Abbe, A. Montanari, "Asymptotic Mutual Information for the Two-Groups Stochastic Block Model" (July 30, 2015); journal version "Asymptotic mutual information for the balanced binary stochastic block model", Information and Inference 6 (2017) 125-170. arXiv:1507.08685v1; doi:10.1093/imaiai/iaw017.
\end{thebibliography}
\space%

\end{document}
